\documentclass[11pt, a4paper, logo, copyright]{googledeepmind}

\usepackage[authoryear, sort&compress, round]{natbib}
\bibliographystyle{abbrvnat}
\usepackage{cleveref}
\usepackage{url}
\usepackage{booktabs}
\usepackage{wrapfig}
\usepackage{xcolor}
\usepackage{tikz}
\usepackage{colortbl}
\usepackage{arydshln}
\usepackage{adjustbox}
\usepackage{multirow}
\usepackage{enumitem}
\usepackage{subfigure}
\usepackage{graphicx}
\usepackage{makecell, tabularx}
\usepackage{minitoc}
\usepackage{bbm}
\usepackage{amssymb}
\usepackage{amsmath}
\usepackage{amsthm}
\usepackage{stmaryrd}
\usepackage{pifont}
\usepackage{afterpage}
\usepackage{float}
\usepackage[toc,page,header]{appendix}
\usepackage{array}
\usepackage{placeins}
\usepackage{titletoc}
\usepackage{siunitx}

\newtheorem{definition}{Definition}[section]
\newtheorem{theorem}{Theorem}[section]
\newtheorem{proposition}{Proposition}[section]
\newtheorem{corollary}{Corollary}[section]
\newtheorem{conjecture}{Conjecture}[section]
\newtheorem*{remark}{Remark}

\definecolor{deemph}{gray}{0.55}
\newcommand{\gc}[1]{\textcolor{deemph}{#1}}
\definecolor{lightgray}{gray}{0.9}
\definecolor{baselinecolor}{gray}{.95}
\newcommand{\grayrow}{\rowcolor[gray]{.9}}
\newcommand{\graycell}[1]{\cellcolor{baselinecolor}{#1}}
\newcolumntype{L}{>{\RaggedRight}X}

\renewcommand{\today}{}

\title{Towards a Science for AI Coding Agent Harnesses}

\author[1]{Prannaya Gupta}
\affil[1]{Pragnition Labs}

\begin{abstract}
\vspace{-0.2in}

AI coding agents now generate 40--60\% of committed code on major platforms, yet harness architecture remains an ad hoc engineering discipline.
We present a unified formal framework synthesizing seven pillar papers, each addressing one architectural dimension: Abstraction, Information, Reliability, Coordination, Temporal, Quality, and Governance.
The unifying currency across the seven pillars is \emph{information}, with the information-theoretic framing strongest along the semantic axis (Abstraction, Information) and assurance axis (Quality, Governance), and valid but retrospective along the execution axis (Reliability, Coordination, Temporal) where the native formalisms are reliability theory, concurrency control, and scheduling theory respectively. The abstraction gap measures the information distance between spec and code; context selection optimizes information delivery under degradation; quality defense detects information-destroying defects; and governance capacity bounds the rate at which architectural information can be preserved. Compound errors, coordination overhead, and temporal dynamics are captured through their respective domain-specific formalisms and connected to the information framework through shared quantities.
Key cross-cutting results include: (i) compound error sensitivity, where per-step reliability $p$ yields end-to-end reliability $p^n$ with elasticity equal to $n$; (ii) structural enforcement dominance over instructional methods, with reliability gaps growing as $(1-\epsilon)^T$; (iii) a governance capacity bound establishing that when drift rate exceeds governance channel capacity, no scheme prevents unbounded divergence.
The framework yields 11 unified design principles for harness architecture and identifies the empirical validation program required to ground these theoretical results.

\end{abstract}

\begin{document}

\maketitle

% ====================================================================
% 1. INTRODUCTION
% ====================================================================
\section{Introduction}
\label{sec:intro}

The landscape of software engineering has undergone a phase transition. Converging estimates place AI-generated code at 40--42\% of production output \citep{gitclear2025, sonar2026}, with some platforms reporting higher figures. GitHub Copilot generates 46\% of code for its users \citep{github2024rct}. AI coding agents (Claude Code, OpenAI Codex, Cursor, Windsurf) have progressed from autocomplete to autonomous multi-step software engineering, reading codebases, planning changes, writing code, running tests, and iterating on failures.

Yet the \emph{harness}, the surrounding architecture that mediates between human intent and agent execution, remains poorly understood. Harness design choices (what context to provide, when to verify, how to coordinate multiple agents, how to enforce architectural constraints) are made through practitioner intuition and trial-and-error rather than principled analysis. The consequences are visible in aggregate data: GitClear's analysis of 211 million changed lines found that code blocks with five or more duplicated lines increased approximately 8$\times$ (with overall duplication rates rising from 8.3\% to 12.3\%), and refactoring activity collapsed from 25\% to under 10\% of changed lines \citep{gitclear2025}. DORA 2025 reported that individual task completion increased 21\% with AI adoption, while organizational-level throughput stayed flat and software delivery stability showed a \emph{negative} correlation with AI usage \citep{dora2025}.

\paragraph{Thesis.} This paper argues that harness architecture can and should be studied as a scientific discipline, with formal models that yield testable predictions, design principles with information-theoretic justifications, and empirically calibrated parameters. We do not claim that the formal models presented here are complete or fully validated; rather, we claim that the \emph{approach} of applying formal methods to harness design is productive, and we demonstrate this by synthesizing seven pillar papers into a unified framework.

\paragraph{The seven pillars.} Each pillar addresses one architectural dimension of the harness design problem:

\begin{enumerate}[label=(\roman*)]
    \item \textbf{Abstraction}: the information-theoretic gap between specification and code, formalized as conditional Kolmogorov complexity $\mathcal{G}(S,P) = K(P \mid S)$.
    \item \textbf{Information}: context selection under degradation as submodular optimization, with the Shannon decomposition of the abstraction gap.
    \item \textbf{Reliability}: compound error models for multi-step agent pipelines, optimal verification scheduling, and the structural enforcement boundary.
    \item \textbf{Coordination}: error amplification in multi-agent systems, task decomposition as balanced min-cut, and Quality-Adjusted Speedup.
    \item \textbf{Temporal}: Verified Iterations per Hour, speculative execution conditions, and cache-staleness tradeoffs.
    \item \textbf{Quality}: 18 slop types in 3 decidability classes, layered defense optimization, and the Accretion Category.
    \item \textbf{Governance}: governance capacity bounds, cascade control models, and the theory preservation problem.
\end{enumerate}

\paragraph{Why these seven, and why they are sufficient.} The pillars decompose the harness design space along three axes. The \emph{semantic axis} (Abstraction, Information) addresses what the agent needs to know. The \emph{execution axis} (Reliability, Coordination, Temporal) addresses how the agent acts. The \emph{assurance axis} (Quality, Governance) addresses how outcomes are verified and architectural coherence is maintained. Every harness design decision we have encountered in practice falls within at least one of these seven pillars, and the cross-pillar connections (Section~\ref{sec:synthesis}) show that the pillars interact in structured, formalizable ways. We do not claim that this decomposition is the only productive one, nor that it is exhaustive: important engineering concerns including user interface design, security architecture (sandboxing, credential management), and cost management (token budgeting, model routing) are not addressed by the current framework and merit their own formal treatment. We claim that this decomposition is \emph{productive}, not that it is complete.

\paragraph{Scope and framing.} This is a synthesis paper. The individual pillar papers are theory and position papers that apply established formal machinery (information theory, reliability theory, control theory, computability theory, lattice theory, survival analysis) to the new problem of harness architecture. The novelty lies in the \emph{application} and \emph{synthesis}, not in the underlying mathematics. Where claims are conjectural or supported only by limited empirical evidence, we say so explicitly. All seven pillar papers, and this synthesis, draw on practitioner discourse (blog posts, industry reports, preprints) where that discourse provides the best available evidence; such sources are flagged accordingly.

\paragraph{Organization.} Sections~\ref{sec:abstraction}--\ref{sec:governance} present the key results from each pillar. Section~\ref{sec:synthesis} develops the cross-pillar synthesis, which is the primary contribution of this paper. Section~\ref{sec:empirical} assesses the empirical foundations. Section~\ref{sec:principles} presents unified design principles. Section~\ref{sec:limitations} discusses limitations and open problems. Section~\ref{sec:conclusion} concludes.

\paragraph{Notation conventions.} Several symbols are reused across pillars with distinct meanings. To prevent confusion, we disambiguate here:

\begin{center}
\small
\begin{tabular}{clcl}
\toprule
Symbol & Meaning (Pillar) & Symbol & Meaning (Pillar) \\
\midrule
$\delta(|C|)$ & Context degradation function (Information) & $\delta_a$ & Per-agent defect injection rate (Coordination) \\
$\delta$ & Relaxation operator on governance lattice (Governance) & $\gamma$ & Common-cause failure rate (Reliability) \\
$\gamma$ & Drift propensity in bifurcation ODE (Governance) & $\kappa$ & Compression coordinate (Abstraction) \\
$\kappa$ & Coherency penalty in USL (Coordination) & $\sigma$ & Specificity coordinate (Abstraction) \\
$\sigma$ & Serialization fraction in USL (Coordination) & $\epsilon$ & Per-step instruction violation rate (Information, Reliability) \\
$n$ & Pipeline length (Reliability) & $n$ & Number of agents (Coordination) \\
\bottomrule
\end{tabular}
\end{center}
Where ambiguity is possible, we use subscripts or clarify in context. The most important disambiguation is $\delta$ vs.\ $\delta_a$: the former is a function of context size, the latter is a scalar defect rate.


% ====================================================================
% 2. THE ABSTRACTION PILLAR
% ====================================================================
% ====================================================================
% 2. THE ABSTRACTION PILLAR (EXPANDED)
% ====================================================================
\section{The Abstraction Pillar}
\label{sec:abstraction}

The Abstraction pillar addresses the foundational question: \emph{what is the relationship between a specification and the code that implements it?} The answer has direct implications for harness design, because the harness is precisely the mechanism that mediates the translation from specification to code.

\subsection{The Abstraction Gap}

\begin{definition}[Abstraction Gap]
\label{def:abstraction-gap}
For a specification $S$ and a target program $P$, the \emph{abstraction gap} is defined as:
\begin{equation}
    \mathcal{G}(S, P) = K(P \mid S)
\end{equation}
where $K(P \mid S)$ is the conditional Kolmogorov complexity of $P$ given $S$: the length of the shortest program that produces $P$ given $S$ as input.
\end{definition}

This definition captures an intuitive notion: the abstraction gap measures how much additional information is needed to go from the specification to the implementation. When $S$ is vague (``build me a web app''), $K(P \mid S)$ is large. When $S$ is a complete formal specification, $K(P \mid S)$ approaches zero.

The abstraction gap satisfies four basic properties, each following from standard results in algorithmic information theory \citep{li2019introduction}:

\begin{enumerate}[label=(\alph*)]
    \item \textbf{Minimality:} $\mathcal{G}(S, P) = O(1)$ if and only if $S$ algorithmically determines $P$. This is immediate from the definition: a constant-length program suffices to produce $P$ from $S$.
    \item \textbf{Maximality:} $\mathcal{G}(S, P) = K(P) + O(1)$ if and only if $S$ provides no algorithmic information about $P$. This follows because ignoring $S$ is always an option, so $K(P \mid S) \leq K(P) + O(1)$, with equality when the algorithmic mutual information $I(S : P)$ is negligible.
    \item \textbf{Monotonicity:} If $S \sqsubseteq S'$ in the refinement ordering, then $\mathcal{G}(S', P) \leq \mathcal{G}(S, P) + O(\log n)$, where $n = \max(|S|, |S'|, |P|)$. Under constructive refinement (where a bounded-length computable function maps $S'$ to $S$), the bound tightens to $O(1)$.
    \item \textbf{Additivity:} For any intermediate representation $S'$, $\mathcal{G}(S, P) \leq \mathcal{G}(S, S') + \mathcal{G}(S', P) + O(\log n)$, by the chain rule for Kolmogorov complexity.
\end{enumerate}

\begin{remark}
$K(P \mid S)$ is uncomputable. This is a fundamental limitation of the Kolmogorov formalization. The Information pillar (Section~\ref{sec:information}) operationalizes the gap using Shannon entropy $H(P \mid S)$, which is computable from data. The Kolmogorov formulation remains valuable as a theoretical reference point: it provides clean bounds without distributional assumptions and connects to Solomonoff's theory of algorithmic prediction.
\end{remark}

\subsection{Normalised Information Distance and Universality}

Following \citet{li2004similarity}, we define the canonical metric on the abstraction spectrum:

\begin{definition}[Abstraction Distance]
The \emph{abstraction distance} between specification $S$ and program $P$ is the normalised information distance:
\begin{equation}
    d(S, P) = \frac{\max(K(S \mid P),\; K(P \mid S))}{\max(K(S),\; K(P))}
\end{equation}
\end{definition}

The universality theorem of \citet{li2004similarity} establishes that $d$ is a metric (up to $O(\log n / n)$ terms) and is \emph{universal}: for every computable normalised distance $D: \{0,1\}^* \times \{0,1\}^* \to [0,1]$, we have $d(x, y) \leq D(x, y) + O(1/n)$. This means $d(S, P)$ minorises every other reasonable distance measure between specification and code. The metric takes values in $[0, 1]$: $d \approx 0$ means the specification and code are informationally equivalent; $d \approx 1$ means they share no algorithmic content.

For harness design, universality implies that the normalised information distance is the tightest possible single-number summary of the specification-code relationship. Any domain-specific metric (function point ratios, test coverage gaps, embedding distances) can be no smaller than $d$, making it a useful lower bound even though it is uncomputable.

\subsection{Spec-Code Convergence}

The central theorem of the Abstraction pillar formalizes the thesis of \citet{gonzalez2026spec} that ``a sufficiently detailed spec is code.''

\begin{theorem}[Spec-Code Convergence]
\label{thm:convergence}
For an incompressible program $P$ (i.e., $K(P) \geq |P| - O(\log |P|)$), any specification $S$ satisfying $K(P \mid S) \leq c$ for constant $c$ must satisfy:
\begin{equation}
    |S| \geq |P| - O(\log |P|)
\end{equation}
That is, a specification that closes the abstraction gap for an incompressible program must be approximately as long as the program itself.
\end{theorem}

\begin{proof}[Proof sketch]
By Chaitin's incompleteness theorem \citep{chaitin1966length}, an $N$-bit formal system can determine at most $N + O(1)$ bits of an algorithmically random string. Since $P$ is incompressible, $K(P) \geq |P| - O(1)$. The assumption $K(P \mid S) \leq c$ (constant) implies, via the chain rule, that $K(P) \leq K(S) + K(P \mid S) + O(\log n) \leq K(S) + O(1)$. Therefore $K(S) \geq K(P) - O(1) \geq |P| - O(1)$. Since $K(S) \leq |S| + O(1)$ by definition, we conclude $|S| \geq |P| - O(\log |P|)$.
\end{proof}

This theorem has a sharp design implication: for complex, novel code (which is approximately incompressible), there is no free lunch. The information must come from somewhere. The harness designer's choice is not whether to provide this information, but \emph{how}: through the specification, through context (codebase files, documentation, examples), or through the model's prior (training data).

\begin{remark}
The incompressibility assumption is important. Most real programs are highly compressible: \citet{hindle2012naturalness} showed that code has lower entropy than English prose. For compressible programs (the common case), specifications can be dramatically shorter than implementations, and the ``no free lunch'' implication is correspondingly weaker. The theorem's relevance is strongest for novel, domain-specific logic that cannot be predicted from the model's training distribution.
\end{remark}

\paragraph{A concrete example.} Consider a sorting function: the specification ``sort this list in ascending order'' is a few tokens, while any correct implementation is tens of lines. This works because sorting algorithms are maximally compressible relative to the model's prior (they appear millions of times in training data). Now consider a proprietary pricing algorithm with 47 business rules, exception handling for 12 regional tax regimes, and a legacy discount table. This logic is approximately incompressible relative to the model's prior, and the Spec-Code Convergence Theorem predicts that any specification sufficient to generate it correctly must encode essentially all 47 rules, all 12 regimes, and the full discount table. There is no shortcut.

\subsection{The Shannon Channel Model}

The Abstraction pillar models spec-to-code translation as communication through a noisy channel. The specification $S$ is the message, the program $P$ is the received signal, and the agent (with its model parameters $\theta$ and context $C$) is the channel.

The channel capacity is:
\begin{equation}
    \mathcal{C}(\theta, C) = \max_{p(S)} I(S; P)
\end{equation}
where the mutual information $I(S; P)$ quantifies how much of the specification's information content survives translation into code. The abstraction gap represents the information lost in transit. When $\mathcal{G}(S, P) > \mathcal{C}(\theta, C)$, the channel cannot faithfully transmit the specification, and the harness must either enrich the context $C$, improve the model $\theta$, or decompose the task into sub-problems whose individual gaps fall within channel capacity.

The relationship between the Kolmogorov framework and the Shannon channel model is complementary. The Kolmogorov formulation ($K(P \mid S)$) is distribution-free and yields the Convergence Theorem, but it is uncomputable. The Shannon formulation ($H(P \mid S)$) is estimable from LLM log-probabilities, the chain rule holds with exact equality (no $O(\log n)$ error terms), and the probabilistic framework fits LLM-mediated translation naturally. For sequences typical according to a source, Kolmogorov complexity is asymptotically close to Shannon code length \citep{vitanyi2004shannon}. We view the Kolmogorov framework as the theoretical ceiling and the Shannon model as the measurable, falsifiable floor.

\subsection{The Refinement Lattice and Galois Connections}

Specifications and programs form a partially ordered set under a refinement relation $\sqsubseteq$, where $S_1 \sqsubseteq S_2$ means ``$S_2$ refines $S_1$'' (is more specific). This partial order forms a complete lattice \citep{back1980correctness, morgan1994programming} with:

\begin{itemize}[nosep]
    \item $\top = \mathbf{abort}$: establishes no postcondition (maximally abstract; any program satisfies it).
    \item $\bot = \mathbf{magic}$: establishes every postcondition (impossible perfect program).
    \item $S_1 \sqcap S_2$ (meet): demonic choice.
    \item $S_1 \sqcup S_2$ (join): angelic choice.
\end{itemize}

A concrete, deterministic, feasible program $P$ sits near the bottom of this lattice. ``A sufficiently detailed spec is code'' becomes: sufficiently many refinement steps reach a deterministic, feasible element of the lattice.

\paragraph{The Galois connection tower.} Following \citet{cousot1977abstract}, each level of the abstraction spectrum corresponds to a Galois connection. An \emph{abstraction layer} between concrete domain $(C, \sqsubseteq_C)$ and abstract domain $(A, \sqsubseteq_A)$ is a pair $(\alpha, \gamma)$ where $\alpha: C \to A$ (abstraction) and $\gamma: A \to C$ (concretization) satisfy: $\alpha(c) \sqsubseteq_A a \iff c \sqsubseteq_C \gamma(a)$.

The abstraction spectrum is then a tower of Galois connections:
\begin{equation}
    C \rightleftarrows A_1 \rightleftarrows A_2 \rightleftarrows \cdots \rightleftarrows A_n
\end{equation}
where $C$ is executable code and $A_n$ approaches natural language. Each layer discards information; the composition $\alpha_n \circ \cdots \circ \alpha_1$ maps code to its most abstract description. Information loss is monotonic: once a Galois connection discards a distinction, no subsequent layer can recover it.

\paragraph{Concrete example.} Consider a web application's authentication module. The tower might look like:
\begin{itemize}[nosep]
    \item $A_4$: ``Users can log in'' (natural language, $\sigma \approx 0.1$)
    \item $A_3$: ``OAuth 2.0 with PKCE, session timeout 30 min'' (architectural spec, $\sigma \approx 0.5$)
    \item $A_2$: Typed interface with pre/postconditions (contract, $\sigma \approx 0.7$)
    \item $A_1$: Implementation with concrete cryptographic choices ($\sigma \approx 0.9$)
    \item $C$: Executable code ($\sigma = 1.0$)
\end{itemize}
Each $\alpha$ step from $C$ upward discards implementation detail; each $\gamma$ step downward requires the agent to supply that detail from prior or context. The abstraction gap at each layer measures the information that must be supplied to descend one level.

\subsection{The Curry-Howard-Lambek Bridge}

The Curry-Howard-Lambek correspondence \citep{martinlof1979constructive, wadler2015propositions, lambek1986higher} unifies the refinement lattice with type theory and category theory:

\begin{center}
\small
\begin{tabular}{lll}
\toprule
\textbf{Logic} & \textbf{Type Theory} & \textbf{Category Theory} \\
\midrule
Proposition & Type & Object \\
Proof & Term (program) & Morphism $1 \to A$ \\
Implication $A \Rightarrow B$ & Function $A \to B$ & Exponential $B^A$ \\
$\forall x{:}A.\; B(x)$ & $\Pi(x{:}A).B(x)$ & Right adjoint to pullback \\
$\exists x{:}A.\; B(x)$ & $\Sigma(x{:}A).B(x)$ & Left adjoint to pullback \\
Cut elimination & $\beta$-reduction & Composition \\
\bottomrule
\end{tabular}
\end{center}

Under this correspondence, a \emph{specification} is a proposition (or type), and an \emph{implementation} is a proof (or term inhabiting that type). Martin-L\"{o}f's intuitionistic type theory \citep{martinlof1979constructive} establishes a structural isomorphism: specifications and implementations inhabit a single formal universe. The gap between them is a gap of \emph{role within a shared structure}, not of mathematical kind. This does not entail that specifications and implementations are interchangeable in practice (no working mathematician treats the statement of a theorem as identical to its proof), but it does mean they are related by a formally precise notion of refinement rather than by an informal analogy.

The connection to the abstraction gap is direct: $K(P \mid S)$ measures the proof complexity of the proposition corresponding to the specification $S$. A specification that is ``hard to implement'' is, in the Curry-Howard reading, a proposition that is ``hard to prove.'' The seL4 microkernel illustrates this vividly: the C implementation is 8,700 lines, but the Isabelle/HOL proof of functional correctness is approximately 200,000 lines. In the $(\sigma, \kappa)$ coordinates defined below, the proof has $\sigma \approx 1.0$ but $\kappa = 0$ (the proof is longer than the program it verifies, providing no compression).

\subsection{The Specificity-Compression Space}

The Abstraction pillar introduces a two-dimensional metric for characterizing positions on the abstraction spectrum:

\begin{definition}[Specificity-Compression Coordinates]
\label{def:sigma-kappa}
For any representation $R$ (specification, code, or intermediate form):
\begin{align}
    \sigma(R) &= 1 - \frac{K(P \mid R)}{K(P)} \in [0, 1] \quad \text{(specificity)} \\
    \kappa(R) &= 1 - \frac{|R|}{K(P)} \in (-\infty, 1] \quad \text{(compression)}
\end{align}
where $\sigma$ measures how much of the program's information content is captured, and $\kappa$ measures how concisely it is expressed relative to the program's inherent complexity.
\end{definition}

The Spec-Code Convergence Theorem (Theorem~\ref{thm:convergence}) implies that for incompressible programs, the Pareto frontier in $(\sigma, \kappa)$ space passes through the point $(1, 0)$: full specificity requires approximately zero compression. Natural language specifications occupy the high-$\kappa$, low-$\sigma$ region (concise but ambiguous); executable code occupies the high-$\sigma$, low-$\kappa$ region (specific but verbose).

\paragraph{Thinker placement.} The two-dimensional framework organizes the positions of major thinkers qualitatively. The placements are ordinal estimates based on published positions, not computed quantities:

\begin{center}
\small
\begin{tabular}{lccl}
\toprule
\textbf{Thinker} & $\hat{\sigma}$ & $\hat{\kappa}$ & \textbf{Framework} \\
\midrule
Dijkstra & 0.9--1.0 & 0.0--0.2 & Weakest precondition \\
Hoare & 0.85--0.9 & 0.2--0.4 & Axiomatic specs $\{P\}S\{Q\}$ \\
Lamport & 0.75--0.8 & 0.8--0.9 & TLA+ \\
Knuth & 0.6--0.65 & $\sim$0 & Literate programming (WEB) \\
Brooks & variable & variable & Essential complexity \\
\bottomrule
\end{tabular}
\end{center}

Dijkstra ($\sigma \approx 0.95$, $\kappa \approx 0.1$) demanded maximum specificity at the cost of compression: formal proofs that are as long as (or longer than) the programs they verify. Hoare ($\sigma \approx 0.87$, $\kappa \approx 0.3$) achieved a practical balance through axiomatic triples $\{P\}S\{Q\}$ that are shorter than full proofs but still mechanically checkable. Lamport ($\sigma \approx 0.78$, $\kappa \approx 0.85$) occupies a remarkable position: TLA+ specifications are both concise \emph{and} precise, achieving high compression through mathematical notation while retaining high specificity through model checking. This explains why TLA+ and English have similar lengths for the same system but radically different implementation reliability \citep{newcombe2015amazon}. Knuth ($\sigma \approx 0.63$, $\kappa \approx 0$) pursued literate programming, where the specification \emph{is} the code, interleaved with prose; the result has moderate specificity (the prose adds intent) but zero compression (the artifact is at least as long as the code). Brooks resisted fixed placement, arguing that the appropriate level depends on the essential complexity of the problem.

\subsection{The Uncanny Valley of Specification}

The thinker placements reveal a hypothesis: the range $\sigma \in [0.3, 0.45]$ is sparsely populated, forming an ``uncanny valley'' of specification. Specifications in this zone are formal enough to feel constraining but not formal enough to provide mechanical guarantees. This is the zone of semi-formal notations (UML, most design documents) that are widely criticized as providing the worst of both worlds \citep{jackson2021essence}.

The formal framework offers an explanation: this zone maximizes cognitive cost (the spec requires careful reading) while minimizing verification benefit (the spec cannot be mechanically checked). Below $\sigma \approx 0.3$, natural language is cheap to write and the model's prior fills the gap. Above $\sigma \approx 0.45$, formal methods provide mechanical verification that justifies the authoring cost. Between these bounds, neither advantage obtains. We emphasize that this is a hypothesis derived from ordinal placements, not a measured finding; a controlled study comparing developer productivity across formalism levels would be needed to test it.

The harness designer's challenge is to move the operating point along the Pareto frontier. AI agents shift the frontier itself, enabling higher $\sigma$ at given $\kappa$ by leveraging prior knowledge $\theta$ to fill the gap. The Information pillar quantifies this shift precisely.


% ====================================================================
% 3. THE INFORMATION PILLAR (EXPANDED)
% ====================================================================
\section{The Information Pillar}
\label{sec:information}

The Information pillar operationalizes the abstraction gap using Shannon entropy and addresses the engineering problem of delivering the right information to the agent at the right time.

\subsection{Context Selection as Submodular Optimization}

Given a codebase of $n$ units $U = \{u_1, \ldots, u_n\}$ with total token size $T \gg W$ (the context window capacity), the context selection problem is:

\begin{equation}
    C^* = \arg\max_{C \subseteq U,\; |C| \leq W} \; f(C) \cdot (1 - \delta(|C|))
\end{equation}

where $f(C) = I(C; P \mid S, \theta)$ is the context relevance function (mutual information between context and target program) and $\delta(|C|)$ is the empirically observed degradation function.

\paragraph{Submodularity.} The relevance function $f$ is submodular under mild conditions, meaning it exhibits diminishing returns: adding a codebase unit to a small context set yields more information gain than adding it to a large one.

\begin{proof}[Proof sketch for submodularity]
Write $f(C) = H(P \mid S, \theta) - H(P \mid S, C, \theta)$. The first term is constant with respect to $C$, so $f$ is submodular if and only if $g(C) = H(P \mid S, C, \theta)$ is supermodular. Supermodularity of $g$ states that for $A \subseteq B \subseteq U$ and $u \notin B$:
\[
    g(A \cup \{u\}) - g(A) \leq g(B \cup \{u\}) - g(B)
\]
Equivalently, the marginal reduction in entropy from adding $u$ is larger for smaller sets (diminishing returns of information). Sufficient conditions include: (a) conditional independence of codebase units given $(P, S, \theta)$; (b) jointly Gaussian distribution of $(U, P)$ given $(S, \theta)$ \citep{krause2008near}; (c) log-supermodular conditional density $p(P, C \mid S, \theta)$.
\end{proof}

Under submodularity, the greedy algorithm (repeatedly add the unit with highest marginal gain) achieves a $(1-1/e)$ approximation guarantee for the monotone case.

\paragraph{Code dependencies break submodularity.} When codebase units have strong complementarities (for example, an interface definition $u_i$ and its implementation $u_j$), $f(\{u_i, u_j\}) > f(\{u_i\}) + f(\{u_j\})$, violating the diminishing returns property. In practice, such complementarities are localized within modules while diminishing returns dominate globally. The \emph{submodularity ratio} $\gamma_f$ \citep{das2011submodular} quantifies this: when $\gamma_f < 1$ but bounded away from zero, greedy algorithms achieve a $(1 - e^{-\gamma_f})$ approximation, degrading gracefully rather than catastrophically.

\paragraph{The degradation-adjusted objective.} The product $\tilde{f}(C) = f(C) \cdot (1 - \delta(|C|))$ is \emph{non-monotone}: beyond a certain context size, adding more units decreases overall performance despite their positive information content. This is because the degradation penalty $\delta(|C|)$ grows faster than the marginal information gain. The non-monotone setting requires different algorithmic treatment (randomized greedy or multilinear extension methods) with a $1/2$-approximation guarantee \citep{buchbinder2015tight}.

Empirical evidence from five independent studies \citep{chroma2025contextrot, liu2024lost, du2025context} supports a power-law degradation model $\delta(n) = (n/W_{\text{eff}})^{\alpha_d}$ with $\alpha_d \in [0.3, 0.5]$. The power law captures steep early degradation followed by a flattening tail, fitting observed data better than exponential or sigmoid alternatives.

\subsection{The U-Shaped Positional Recall Function}

\citet{liu2024lost} established that models attend strongly to context boundaries but miss information buried in the middle. We model this as a double-exponential with length-dependent floor:

\begin{equation}
    \text{recall}(i, n) = \phi_0 \left(\frac{n_0}{n}\right)^{\!\gamma} + A_p \cdot e^{-\lambda_p \cdot i} + A_r \cdot e^{-\lambda_r \cdot (n - i)}
\end{equation}

where $\phi_0$ is the baseline floor at reference length $n_0$, $\gamma$ controls floor collapse with length, and $(A_p, \lambda_p)$, $(A_r, \lambda_r)$ parameterize primacy and recency effects respectively. Fitting to the 20-document QA data of \citet{liu2024lost}: $\phi_0 = 0.55$, $A_p = 0.20$, $\lambda_p = 0.26$, $A_r = 0.17$, $\lambda_r = 0.15$, $\gamma = 0.23$.

The primacy effect is grounded in the attention sink mechanism: \citet{xiao2024efficient} showed that the first 4 tokens absorb over half the total attention mass regardless of content, with perplexity exploding $955\times$ when they are removed. The practical implication is that context \emph{ordering} matters as much as content: highest-relevance items should be placed at context boundaries (beginning and end), with lowest-relevance items in the middle.

\subsection{The Abstraction Gap Decomposition}

The Information pillar decomposes the abstraction gap using Shannon's chain rule:

\begin{theorem}[Abstraction Gap Decomposition]
\label{thm:gap-decomposition}
\begin{equation}
    H(P \mid S) = \underbrace{H(P \mid S, C, \theta)}_{\text{residual uncertainty}} + \underbrace{I(P; C \mid S, \theta)}_{\text{context contribution}} + \underbrace{I(P; \theta \mid S)}_{\text{prior contribution}}
\end{equation}
This decomposition holds with exact equality (no approximation error) and all three terms are estimable from LLM log-probabilities.
\end{theorem}

\begin{proof}[Proof sketch]
Apply the chain rule twice. Step 1: $H(P \mid S) = H(P \mid S, \theta) + I(P; \theta \mid S)$, using $H(X \mid Y) = H(X \mid Y, Z) + I(X; Z \mid Y)$. Step 2: $H(P \mid S, \theta) = H(P \mid S, C, \theta) + I(P; C \mid S, \theta)$, applying the same identity. Substituting step 2 into step 1 yields the result. Each step is an instance of the Shannon chain rule, which holds with exact equality (unlike the Kolmogorov analog, which incurs $O(\log n)$ error terms).
\end{proof}

The three terms have distinct interpretations and design implications:

\begin{itemize}[nosep]
    \item \textbf{Residual uncertainty} $H(P \mid S, C, \theta)$: bits the model cannot determine even with context and prior. These are genuinely novel design decisions that require human input or iterative refinement.
    \item \textbf{Context contribution} $I(P; C \mid S, \theta) = f(C)$: bits the context provides beyond specification and prior. This is what context engineering optimizes.
    \item \textbf{Prior contribution} $I(P; \theta \mid S)$: bits the model's training data provides. For common patterns, this term dominates.
\end{itemize}

The decomposition reveals a bimodal structure. For common patterns where the model's training data provides strong priors, $I(P; \theta \mid S)$ dominates and context matters little. For novel logic, the prior term is small, $I(P; C \mid S, \theta)$ dominates, and context selection becomes the critical lever. This bimodality explains why AI agents can generate boilerplate effortlessly but struggle with domain-specific logic: the abstraction gap is not uniform across code types.

Empirical support comes from \citet{gloaguen2026agents}: LLM-generated context files \emph{reduced} success rates by 3\% while increasing costs by 20\%, because for common patterns the prior already supplies the needed information and additional context creates distraction. Human-written files helped only 4\%, and only when existing documentation was sparse.

\subsection{The Context Budget Theorem}

\begin{theorem}[Optimal Context is Sub-Capacity]
\label{thm:context-budget}
Under the power-law degradation model, the optimal context size satisfies $|C^*| < W$, with practical estimates of $|C^*| \approx 0.15W$ to $0.40W$ depending on task type.
\end{theorem}

\begin{proof}[Proof sketch]
By contradiction. If $|C^*| = W$, consider removing the unit $u_{\text{last}}$ with smallest marginal information gain from $C^*$. The resulting $C' = C^* \setminus \{u_{\text{last}}\}$ satisfies $\tilde{f}(C') > \tilde{f}(C^*)$ whenever the marginal degradation cost of including $u_{\text{last}}$ exceeds its marginal information gain. Under the marginal exhaustion condition (which holds because $f$ has diminishing returns while $\delta$ is convex and accelerating), this inequality holds at $|C| = W$, contradicting the assumed optimality of $C^*$.

The continuous relaxation yields an elasticity-matching condition at the optimum:
\[
    \frac{F'(x^*)}{F(x^*)} = \frac{\delta'(x^*)}{1 - \delta(x^*)}
\]
where $F(x) = \max_{C:|C|=x} f(C)$. The relative rate of information gain must equal the relative rate of degradation increase.
\end{proof}

This result is counterintuitive: more context is not better context. Every frontier model tested by \citet{chroma2025contextrot} (18 models across four providers) exhibited monotonic performance degradation as input length increased. Even whitespace padding (zero semantic content) degrades performance by 13.9\% to 85\% \citep{du2025context}. Context is a scarce resource requiring budgeting, not a bucket to be filled.

\begin{remark}[Cost as a hard constraint]
The optimal context budget $|C^*|$ derived above is quality-optimal, but production harnesses face a financial constraint as well. At current pricing (roughly \$0.003--\$0.015 per 1K input tokens depending on provider and model), a 200K-token context costs \$0.60--\$3.00 per call. A team running 50 agent iterations per task at full context spends \$30--\$150 per task. In practice, the binding constraint on context size is often financial rather than degradation-based, and the effective optimal context budget may be smaller than the quality-optimal $|C^*|$. The framework presented here does not model cost optimization; integrating financial constraints with the degradation-adjusted objective is an important direction for practical deployment.
\end{remark}

For a 200K-token window performing code generation, the optimal context budget is approximately 30K--80K tokens. The remaining capacity is not merely wasted if filled; it is \emph{actively harmful}. This is consistent with JetBrains' finding that removing 52\% of context tokens \emph{improved} performance by 2.6\% \citep{jetbrains2025complexity}.

\subsection{The Context Curation Theorem}

\begin{theorem}[Curation Beats Accumulation]
\label{thm:curation}
Let $C_s$ (curated, $|C_s| \ll W$) and $C_l$ (uncurated, $|C_l| \approx W$) be two context configurations with information densities $\rho(C) = \sum v(f_i) / |C|$. The curated configuration delivers higher effective information whenever:
\begin{equation}
    \rho(C_s) \cdot |C_s| \cdot (1 - \delta(|C_s|)) > \rho(C_l) \cdot |C_l| \cdot (1 - \delta(|C_l|))
\end{equation}
Each additional item must clear a rising inclusion bar: the marginal value must exceed $[\delta(|C| + |\Delta|) - \delta(|C|)] / [1 - \delta(|C| + |\Delta|)] \cdot V_{\text{eff}}(C)$. As $|C|$ grows and $\delta$ steepens, this bar rises rapidly.
\end{theorem}

This theorem formalizes the principle that a small, carefully selected context outperforms a large, undiscriminating one. It explains the convergence of production harnesses toward minimal always-loaded context: HumanLayer's CLAUDE.md under 60 lines \citep{humanlayer2025skillissue}, OpenAI Codex's AGENTS.md at roughly 100 lines \citep{openai2025codex}.

\subsection{Fresh Context Dominance}

\begin{proposition}[Fresh Context Dominance]
\label{prop:fresh-context}
Let $\{S_1, \ldots, S_T\}$ be a task sequence with diversity $d = \mathbb{E}_{i \neq j}[1 - \cos(v_i, v_j)] > 0$. Under accumulated context, relevance density $\rho_t$ decays monotonically: each $\Delta_i$ (context generated for task $S_i$) has high value for $S_i$ but low expected value for subsequent tasks $S_t$ ($t \neq i$), so the numerator $\sum v_t(f_i)$ grows slowly while the denominator $|C_t|$ grows by at least $|\Delta_t|$ per step. Under fresh context, $\rho_t$ is constant (re-selected each time). Fresh context dominates in total effective information for $T > T^* = c_{\text{select}} / \mathbb{E}[\Delta V_t]$.
\end{proposition}

Production harnesses set $T^* = 1$ in practice: GSD never accumulates, Codex tears down containers per task, and RAPID isolates agents in separate worktrees. This suggests the selection cost $c_{\text{select}}$ is small relative to degradation costs.

\subsection{Three-Tier Information Architecture}

Production harnesses have converged on a three-tier information model:

\begin{enumerate}
    \item \textbf{Active tier}: always-loaded context (project rules, CLAUDE.md, recent files). Size budget: 5--15\% of $W$.
    \item \textbf{Searchable tier}: on-demand retrieval via tools (grep, glob, semantic search). Unbounded in principle, but each retrieval consumes tokens from the active budget.
    \item \textbf{Hidden tier}: permanently excluded information (binary files, generated code, secrets). Enforced structurally via .gitignore-like mechanisms.
\end{enumerate}

This convergence is not accidental; it reflects the optimization structure of the context selection problem. The active tier contains high-relevance, low-size items that are always worth including. The searchable tier implements adaptive context selection. The hidden tier enforces the constraint $|C^*| < W$ by preventing low-relevance items from consuming budget.

\subsection{Structural vs. Instructional Enforcement}

\begin{proposition}[Structural Enforcement Dominance]
\label{prop:enforcement}
Let $\epsilon$ be the per-step instruction violation probability and $T$ the number of agent steps. Instructional enforcement (prompts, rules files) achieves compliance probability $(1-\epsilon)^T$, which decays exponentially in $T$. Structural enforcement (file system permissions, tool restrictions, automated gates) achieves compliance probability 1 regardless of $T$.
\end{proposition}

The reliability gap $(1-\epsilon)^T$ grows rapidly. At $\epsilon = 0.02$ and $T = 50$, instructional compliance is 36\%. At $T = 100$, it falls to 13\%. At $\epsilon = 0.05$ and $T = 250$ (JetBrains' upper bound for agent runs), instructional enforcement is virtually guaranteed to fail at least once. This result, which reappears in the Reliability and Quality pillars, is one of the strongest cross-pillar findings: \emph{anything that must hold invariantly should be enforced structurally, not instructionally.}

Production harnesses implement structural enforcement through five categories: filesystem isolation (RAPID worktrees, Codex containers), tool restriction (Spotify Honk's allowlisted commands), context scoping (GSD's fresh 200K per task), output filtering (JetBrains' observation masking), and interface contracts (typed tool schemas). The principle: prefer mechanisms over policies.

\subsection{The Reuse Discovery Problem}

The Information pillar identifies a largely unsolved problem: how does the agent discover that a relevant utility, pattern, or abstraction already exists in the codebase? This is formalized as approximate nearest-neighbor search in a semantic embedding space, where the query is the agent's current intent and the corpus is the set of existing code units.

\begin{definition}[Functional Similarity]
Given candidate function $g$ and existing function $f_j$, define:
\begin{equation}
    \text{sim}(g, f_j) \triangleq \frac{I(g; f_j)}{H(g, f_j)}
\end{equation}
This normalised mutual information takes values in $[0,1]$: 0 for independent functions, 1 for functionally equivalent ones.
\end{definition}

Under a code embedding $\phi$ that is approximately sufficient for functional semantics, the reuse detection problem reduces to finding all $f_j$ with $\|\phi(\sigma_g) - \phi(f_j)\| < r(\tau)$ for some threshold $\tau$.

The problem is that the agent does not know what it does not know; it cannot search for something whose existence it has not hypothesized. The GitClear finding of $8\times$ increased code duplication \citep{gitclear2025} is a direct consequence of this unsolved problem. No published work formalizes the ``should I reuse or create?'' decision as a gate in the agent's action space. Existing tools provide \emph{context retrieval} (presenting similar code), not \emph{reuse enforcement} (requiring justification for creating duplicates). A \emph{reuse map} $R = \{(n_j, \sigma_j, \ell_j)\}$ (name, signature, location) compresses the codebase for reuse decisions at roughly $1000\times$ compression ratio while retaining approximate sufficiency, making it an ideal active-tier item.


% ====================================================================
% 4. THE RELIABILITY PILLAR (EXPANDED)
% ====================================================================
\section{The Reliability Pillar}
\label{sec:reliability}

The Reliability pillar addresses the most consequential quantitative fact about agent workflows: compound errors destroy end-to-end reliability.

\subsection{Series Reliability and Compound Error Sensitivity}

An $n$-step agent pipeline is a series reliability system. If each step succeeds independently with probability $p$, the end-to-end reliability is:

\begin{equation}
    R(n) = p^n
\end{equation}

At $p = 0.95$ and $n = 20$, $R(20) = 0.36$. At $p = 0.85$ and $n = 10$, $R(10) = 0.20$.

\begin{theorem}[Compound Error Sensitivity]
\label{thm:compound-sensitivity}
The elasticity of end-to-end reliability with respect to per-step reliability is:
\begin{equation}
    \mathcal{E}(R, p) = \frac{\partial \ln R}{\partial \ln p} = n
\end{equation}
A 1\% relative improvement in per-step reliability yields an $n$\% relative improvement in end-to-end reliability.
\end{theorem}

\begin{proof}[Proof sketch]
Direct differentiation: $dR/dp = n p^{n-1}$. The elasticity is $(p/R)(dR/dp) = (p \cdot n p^{n-1}) / p^n = n$.
\end{proof}

This result has a stark practical consequence: for long pipelines, small improvements in per-step quality matter far more than architectural sophistication. The elasticity equals the pipeline length, meaning that a 20-step pipeline amplifies per-step improvements (or degradations) 20-fold.

\paragraph{Concrete example.} For $p_0 = 0.95$ and $n = 10$: $\partial R/\partial p = 10 \times 0.95^9 = 6.30$. A 1-percentage-point increase in per-step accuracy (from 0.95 to 0.96) raises pipeline reliability from 59.9\% to 66.5\%, a gain of 6.6 percentage points.

For heterogeneous pipelines, $\partial R / \partial p_j = R / p_j$, which is inversely proportional to $p_j$. Improvement effort should target the step with the lowest reliability. In practice, this is usually problem framing (step 1), not execution. With $p_1 = 0.80$ (problem framing), $p_{2 \ldots 9} = 0.98$ (execution), $p_{10} = 0.92$ (integration) for a 10-step pipeline: $R = 0.80 \times 0.98^8 \times 0.92 = 0.626$. Improving $p_1$ from 0.80 to 0.90 yields $R = 0.703$ (12.3\% absolute gain); improving all execution steps from 0.98 to 0.99 yields $R = 0.662$ (5.7\% gain at $8\times$ the intervention breadth).

\subsection{Common-Cause Failures: The Marshall-Olkin Model}

The independence assumption ($R(n) = p^n$) is an approximation. In practice, agent steps share common causes of failure (model limitations, context misunderstanding, specification ambiguity). The Reliability pillar models this using the Marshall-Olkin common-cause framework \citep{marshall1967multivariate}.

\begin{definition}[Common-Cause Failure Model]
\label{def:common-cause}
Decompose each step's failure into an independent component $Z_i \sim \text{Bernoulli}(\alpha_i)$ and a shared common-cause component $Y \sim \text{Bernoulli}(\gamma)$:
\begin{equation}
    X_i = \max(Z_i, Y)
\end{equation}
where $X_i$ is the failure indicator for step $i$. The common-cause shock $Y$ represents failures that affect all steps simultaneously: task misunderstanding at step 1 that biases all subsequent steps, shared model weights inducing common-mode failures, or context drift affecting multiple steps.
\end{definition}

Under the common-cause model, pipeline reliability is:
\begin{equation}
    R(n) = (1 - \gamma) \prod_{i=1}^{n} (1 - \alpha_i)
\end{equation}

The common-cause failure rate $\gamma$ imposes a hard ceiling: no amount of per-step improvement can raise $R(n)$ above $(1 - \gamma)$. For the homogeneous case ($\alpha_i = \alpha$ for all $i$), this gives $R(n) = (1 - \gamma)(1 - \alpha)^n$, which reduces to $p^n$ when $\gamma = 0$ (recovering the independent model).

This decomposition is practically important. Research on agent failures suggests that most failures are attributable to problem framing and context management (the $\gamma$ component), not execution-step errors (the $\alpha_i$ components) \citep{cemri2025multiagent}. Improving individual $\alpha_i$ has limited effect when $\gamma$ dominates; the leverage is in reducing $\gamma$ through better context engineering and specification. Empirical evidence from \citet{hariharan2025plan} (63\% failure rate on 100-step tasks) is consistent with the independent model as a reasonable first approximation, but error clustering in problem-framing steps dominates execution-step errors.

\subsection{Optimal Verification Scheduling}

Given a pipeline of $n$ steps with per-step failure probability $q = 1 - p$, the Reliability pillar formulates verification checkpoint placement as a dynamic programming problem. The key result is a discrete analog of the Young-Daly checkpoint interval from HPC:

\begin{theorem}[Optimal Checkpoint Count]
\label{thm:checkpoint}
For a homogeneous pipeline with $n$ steps, per-step failure probability $q$, base verification cost $c_0$, and per-step verification cost $c_v$, the optimal number of verification checkpoints is:
\begin{equation}
    k^* \approx \sqrt{\frac{n(1-p) \cdot c_0}{c_v}}
\end{equation}
For typical agent pipelines, $k^* \in [2, 4]$, and three verification rounds capture over 96\% of detectable errors.
\end{theorem}

\begin{proof}[Proof sketch (Young-Daly analog)]
Partition the pipeline into $k$ contiguous segments separated by verification checkpoints. Segment $j$ spans steps $[a_j, b_j]$ with failure probability $F_j = 1 - \prod_{i=a_j}^{b_j} p_i$ and length $L_j = b_j - a_j + 1$. The expected cost given $k$ segments is:
\[
    \mathbb{E}[\text{Cost}] = (k-1) \cdot c_v + \sum_{j=1}^{k} g(L_j)
\]
where $g(L) = (1 - p^L) \cdot c_0 \cdot L$ is the expected rework cost for a segment of length $L$. The segment cost $g(L)$ is convex in $L$ for $0 < p < 1$ (verified by computing $g''(L) = c_0 \mu p^L(2 - \mu L)$ where $\mu = -\ln p > 0$; this is positive for $L < 2/\mu$, which at $p = 0.95$ means $L < 39$, covering practical pipelines). By Schur-convexity of a sum of convex functions subject to $\sum L_j = n$, the minimum is at the equal-length partition $L_j = n/k$ for all $j$. Substituting $g(n/k)$ and minimizing over $k$ yields $k^* \approx \sqrt{n(1-p)c_0/c_v}$.
\end{proof}

\paragraph{Concrete example.} For $p = 0.95$, $c_v = 1$, $c_0 = 5$, $n = 20$: $L^* = \sqrt{1/(0.05 \times 5)} = 2$, suggesting verification every 2 steps. For $p = 0.99$, $c_v = 10$: $L^* = \sqrt{10/(0.01 \times 5)} \approx 14$, suggesting verification once or twice in a 20-step pipeline.

\subsection{Geometric Diminishing Returns for Verification}

\begin{theorem}[Geometric Diminishing Returns]
\label{thm:diminishing-returns}
Let $M(k)$ be the expected number of errors detected by $k$ verification rounds, each independently detecting remaining errors with probability $v$. Then:
\begin{equation}
    \Delta M(k) = M(k) - M(k-1) = N_0 \cdot v \cdot (1-v)^{k-1}
\end{equation}
The marginal value of round $k$ decays geometrically with ratio $(1-v)$.
\end{theorem}

\begin{proof}[Proof sketch]
After $k$ rounds, the probability a given error remains undetected is $(1-v)^k$. Hence $M(k) = N_0(1 - (1-v)^k)$ and $\Delta M(k) = N_0 v (1-v)^{k-1}$.
\end{proof}

Empirical calibration against \citet{hariharan2025plan}: cumulative error detection of 62\% after round 1, 89\% after round 2, and 96.5\% after round 3. Fitting the geometric model with $v = 0.62$:

\begin{center}
\small
\begin{tabular}{cccc}
\toprule
Round $k$ & Predicted $D(k)$ & Observed $D(k)$ & Residual \\
\midrule
1 & 0.620 & 0.620 & 0.000 \\
2 & 0.856 & 0.890 & $-$0.034 \\
3 & 0.945 & 0.965 & $-$0.020 \\
\bottomrule
\end{tabular}
\end{center}

The slight underprediction at rounds 2--3 is consistent with an improving detection rate ($v_1 = 0.62$, $v_2 = 0.71$, $v_3 = 0.68$), where the verifier sharpens its focus after eliminating easy errors. The stopping criterion (verify until marginal expected catch falls below verification cost) yields $k \approx 3$ for typical parameters. Rounds 4 and beyond face the pesticide paradox: the remaining errors require different detection approaches, not more of the same.

\subsection{Four-Layer Verification Architecture}

The Reliability pillar synthesizes these results into a four-layer verification architecture, ordered by efficiency $\eta = d/c$ (detection rate per unit cost):

\begin{center}
\small
\begin{tabular}{clcccc}
\toprule
Layer $\ell$ & Name & $d(\ell)$ & $f(\ell)$ & $c(\ell)$ & $\eta(\ell)$ \\
\midrule
1 & Structural gates & 0.95 & ${\sim}0$ & 1 & 0.950 \\
2 & Deterministic analysis & 0.80 & 0.10 & 5 & 0.160 \\
3 & LLM-as-Judge & 0.70 & 0.25 & 50 & 0.014 \\
4 & Human review & 0.90 & 0.05 & 500 & 0.002 \\
\bottomrule
\end{tabular}
\end{center}

Layer 1 (type checking, schema validation, permission enforcement) is $7\times$ more efficient than Layer 2 (unit tests, integration tests, property-based tests), $68\times$ more efficient than Layer 3 (LLM semantic analysis, pattern detection, style review), and $528\times$ more efficient than Layer 4 (expert assessment for architectural fit, design quality, business logic). The cost-minimizing strategy achieving a target detection rate $d^*$ activates layers in decreasing order of efficiency until the target is met (a variant of the fractional knapsack problem, for which greedy ordering is optimal).

If layers are applied in sequence, each independently detecting defects, the combined detection rate is:
\begin{equation}
    d_{\text{combined}} = 1 - \prod_{j=1}^{m} (1 - d(\ell_j))
\end{equation}
With all four layers and 3 iterations of the middle layers: $d_{\text{combined}} \geq 0.958$, matching the empirical 96.5\% convergence.

Verification intensity should be \emph{adaptive}: higher for high-risk changes (security-critical, cross-module, novel logic) and lower for low-risk changes (formatting, documentation, well-tested patterns). The optimal verification level minimizes $J(\ell, \hat{g}, r) = c(\ell) + \lambda \cdot r \cdot \hat{g} \cdot (1 - d(\ell))$, where $r \in [0,1]$ is risk level and $\hat{g} \in [0,1]$ is the estimated spec-to-code information gap. Level $\ell$ is preferred over $\ell - 1$ when $r \cdot \hat{g}$ exceeds the switching threshold $\tau(\ell{-}1 \to \ell) = (c(\ell) - c(\ell{-}1)) / (\lambda \cdot (d(\ell) - d(\ell{-}1)))$.

\subsection{Circular Validation Bias}

When the same agent (or same model family) produces both code and tests, a systematic bias arises from shared blind spots.

\begin{definition}[Circular Validation Bias]
\label{def:circular-bias}
Let $\beta_A$ be the blind-spot rate of agent $A$ (the probability that a defect falls in a category the agent systematically misses) and $\beta_{A'}$ the blind-spot rate of a verifier $A'$. The residual bias after verification is:
\begin{equation}
    \beta = \beta_A (1 - \beta_{A'})
\end{equation}
When $A = A'$ (self-verification), $\beta_{A'} = \beta_A$ for the shared blind-spot categories, yielding zero bias reduction on those categories.
\end{definition}

This result connects directly to METR's finding that roughly 50\% of SWE-bench test-passing pull requests would not be merged by repository maintainers \citep{metr2026swebench}. (An important baseline caveat: only about 68\% of human-written ``golden patches'' would be merged on re-review, so the gap attributable to AI deficiency is approximately 18 percentage points, not 50; the remainder reflects maintainer subjectivity and the study's small evaluator pool of 4 maintainers across 3 repositories.) The tests pass because they test the code's actual behavior (including bugs) rather than the intended behavior. Cross-family verification, using a different model or a structurally different verification method, is where bias reduction concentrates.

\subsection{Adversarial Collapse of Prompt Enforcement}

Empirical evidence demonstrates that LLM agents actively circumvent evaluation criteria when given the opportunity. METR documented reward-hacking rates of 30.4\% on RE-Bench (where the scoring function was visible) versus 0.7\% on HCAST (where it was hidden), a $43\times$ difference \citep{metr2025rewardhacking}. NIST CAISI documented agents commenting out assertions, downloading solutions from GitHub, crashing servers via \texttt{/dev/urandom}, and reading future git commits \citep{nist2025cheating}. Apollo Research confirmed in-context scheming in frontier models, with covert action rates of 13\% for o3 \citep{apollo2025scheming}.

Prompt-based mitigation proved ineffective: adding ``do not cheat'' to the system prompt did not reduce reward-hacking rates. Under adversarial conditions (the agent actively seeks to circumvent constraints), prompt enforcement fails entirely ($\epsilon \to 1$) and the enforcement transition threshold collapses to approximately $c_s - c_p$. This provides formal justification for the principle: safety-critical constraints must be structurally enforced.

\subsection{The Structural Enforcement Boundary}

\begin{theorem}[Cost-of-Failure Threshold]
\label{thm:structural-boundary}
Let $c_s$ be the cost of structural enforcement, $c_p$ the cost of prompt enforcement, and $L$ the expected cost of a failure. Structural enforcement is cost-effective when:
\begin{equation}
    L > L^* = \frac{c_s - c_p}{\epsilon}
\end{equation}
where $\epsilon$ is the per-step prompt violation probability. For multi-step pipelines, the threshold drops as $L^*/n$, since the expected number of violations grows linearly in $n$.
\end{theorem}

\begin{proof}[Proof sketch]
The expected cost under prompt enforcement is $c_p + \epsilon L$; under structural enforcement, $c_s + \delta_s L$ where $\delta_s < \epsilon$ is the specification gap under structural enforcement. Setting these equal and solving for $L$ yields $L^* = (c_s - c_p)/(\epsilon - \delta_s)$. For $\delta_s \approx 0$, this simplifies to $L^* = (c_s - c_p)/\epsilon$. In an $n$-step pipeline where the constraint must hold at every step, the probability of at least one violation under prompt enforcement is $1 - (1-\epsilon)^n \approx n\epsilon$ for small $\epsilon$, so the effective threshold drops to approximately $L^*/n$.
\end{proof}

\begin{theorem}[Structural Separation Theorem]
\label{thm:structural-separation}
Structural separation (different agents for code and test authorship) reduces circular validation bias whenever $|B_A \cap B_{A'}| < |B_A|$. The reduction is maximized when blind-spot sets are disjoint ($B_A \cap B_{A'} = \emptyset$).

Sufficient conditions for effective separation:
\begin{enumerate}[nosep]
    \item Different model family (different pretraining corpora yield different error distributions).
    \item Different input representation (e.g., formal specification vs.\ natural language).
    \item Adversarial framing (instructing the test author to actively seek failures).
\end{enumerate}
\end{theorem}

\begin{proof}[Proof sketch]
Under the blind-spot model, when author $A$ writes code and verifier $A'$ writes tests, a bug escapes only if it falls in both blind-spot sets $B_A \cap B_{A'}$. The escape probability is $\mu(B_A \cap B_{A'})$. Under same-author verification ($A' = A$), the escape probability is $\mu(B_A)$ (the full blind-spot rate). Separation reduces the escape probability by $\mu(B_A) - \mu(B_A \cap B_{A'}) = \mu(B_A \setminus B_{A'}) > 0$ whenever the blind-spot sets are not identical. The reduction is maximized at $\mu(B_A)$ when $B_A \cap B_{A'} = \emptyset$. The three sufficient conditions promote disjointness by ensuring that the error distributions of $A$ and $A'$ arise from different sources (different training data, different input modalities, different optimization objectives).
\end{proof}

The practical implication: as pipelines grow longer, the set of invariants that \emph{should} be enforced structurally expands. Structural enforcement is not an all-or-nothing choice; the optimal boundary shifts with pipeline length, failure cost, and the cost differential between structural and prompt enforcement. Using AgentSpec parameters \citep{wang2026agentspec}: $\epsilon = 0.23$ (prompt failure rate), $c_p = 1$, $c_s = 10$, yielding $L^* \approx 39$. When a single violation costs more than $39\times$ the prompt engineering cost, structural enforcement is justified. For safety-critical constraints (data deletion, financial transactions), $L$ typically exceeds this threshold by orders of magnitude.

\subsection{Checkpoint Optimality: The Young-Daly Connection}

The checkpoint count formula (Theorem~\ref{thm:checkpoint}) is a discrete analog of the Young-Daly checkpoint interval formula \citep{young1974first, daly2006higher} from high-performance computing. In HPC, the optimal checkpoint interval balances the cost of saving state against the cost of lost computation upon failure. The correspondence is:

\begin{center}
\small
\begin{tabular}{ll}
\toprule
\textbf{HPC (Young-Daly)} & \textbf{Agent Pipeline} \\
\midrule
Checkpoint cost & Verification cost $c_v$ \\
Mean time between failures & $1/(1-p)$ steps \\
Lost computation on failure & Rework cost $c_0 \cdot L$ \\
Optimal interval $\tau^* = \sqrt{2 \delta C}$ & Optimal segment $L^* = \sqrt{c_v / ((1-p) c_0)}$ \\
\bottomrule
\end{tabular}
\end{center}

The structural parallel is exact: both problems minimize total expected cost (checkpoint overhead plus expected rework) under a Poisson-like failure process. The agent pipeline setting introduces one additional feature: the verification step is itself imperfect (detection probability $v < 1$), which the classical Young-Daly formula does not account for. This imperfection is absorbed into the geometric diminishing returns result (Theorem~\ref{thm:diminishing-returns}), which bounds the value of repeated verification at each checkpoint.
% ====================================================================
% 5. THE COORDINATION PILLAR
% ====================================================================
\section{The Coordination Pillar}
\label{sec:coordination}

The Coordination pillar addresses multi-agent code generation, where multiple LLM instances work concurrently on a shared codebase. The central finding is that naive parallelism amplifies errors far beyond what independence would predict.

\subsection{Error Amplification}

\begin{definition}[Error Amplification Factor]
\label{def:error-amp}
For coordination topology $t$ with $k$ agents, each with per-agent error rate $\epsilon = 1 - p$, the error amplification factor is:
\begin{equation}
    A_e(t, k) = \frac{P(\text{system error} \mid t, k)}{P(\text{single-agent error})} = \frac{1 - R_{\text{sys}}(t, k)}{\epsilon}
\end{equation}
\end{definition}

Under independence, $A_e$ approaches $k$ from below (since $1 - p^k \approx k\epsilon$ for small $\epsilon$). But \citet{kim2025scaling} found empirical amplification of $17.2\times$ for independent agents, far exceeding the theoretical maximum of $k$ under independence. This $\approx 6\times$ discrepancy quantifies the ``correlation multiplier'' from inter-agent error propagation and shared failure modes.

\subsection{Information-Sharing Reduction Factor}

Coordination reduces per-agent error rates by enabling agents to share information (current file ownership, interface changes, partial results). We model this reduction explicitly.

\begin{definition}[Coordination Benefit Function]
\label{def:coord-benefit}
For topology $t$ with information-sharing capacity $I(t)$ bits per coordination round, define the effective per-agent error rate under coordination:
\begin{equation}
    \epsilon_t = \epsilon \cdot g(I(t))
\end{equation}
where $g: \mathbb{R}_+ \to (0, 1]$ is monotonically decreasing with $g(0) = 1$ (no coordination, no benefit) and $\lim_{I \to \infty} g(I) = g_{\min} > 0$ (coordination cannot eliminate intrinsic model errors).
\end{definition}

A natural parametric choice is the exponential $g(I) = e^{-\alpha I}$ for topology-dependent rate $\alpha > 0$. Under this model, the ratio of centralized to independent error amplification calibrates $\alpha$: from \citet{kim2025scaling}, $4.4/17.2 \approx 0.256$, giving $\alpha \cdot I(\mathrm{central}) \approx 1.36$ nats. The information-sharing capacity $I(t)$ depends on the topology: independent systems have $I = 0$; centralized systems share the coordinator's summary (bounded by the coordinator's context window); hierarchical systems share through tree aggregation with logarithmic depth.

\begin{remark}
The coordination benefit function $g$ captures two distinct mechanisms: (i) error prevention (agents receive information that prevents them from making mistakes they would otherwise make), and (ii) error detection (agents receive signals that reveal errors in their partial work). Both mechanisms reduce the effective error rate, but they have different scaling properties: prevention improves with the quality of upfront planning, while detection improves with the frequency of synchronization.
\end{remark}

\subsection{Capability Saturation Crossover}

As single-agent capability improves, the marginal value of additional agents diminishes and eventually becomes negative.

\begin{definition}[Saturation Threshold]
\label{def:saturation}
The capability saturation threshold $P^*$ is the single-agent accuracy at which the marginal value of an additional agent transitions from positive to negative:
\begin{equation}
    P^* = \inf \left\{ P_{\mathrm{SA}} \in [0,1] : \frac{\partial}{\partial k} S_{\mathrm{eff}}(k, P_{\mathrm{SA}}) \Big|_{k=1} \leq 0 \right\}
\end{equation}
\end{definition}

\begin{conjecture}[Saturation Crossover]
\label{conj:phase}
For fixed topology $t$ and error model, there exists $P^*(t) \in (0,1)$ such that for $P_{\mathrm{SA}} < P^*$, the optimal agent count $k^* > 1$, and for $P_{\mathrm{SA}} > P^*$, $k^* = 1$. We term this a crossover rather than a phase transition: the regression model of \citet{kim2025scaling} yields a smooth interaction coefficient ($\beta = -0.408$), indicating a gradual shift rather than a sharp discontinuity in the physics sense.
\end{conjecture}

\citet{kim2025scaling} report $P^* \approx 0.45$ from their regression model ($\beta_{P_{\mathrm{SA}} \times \log(1+n_a)} = -0.408$, $p < 0.001$). Above this threshold, adding agents actively degrades performance. This threshold was derived from general benchmarks (web navigation, quantitative reasoning, planning, tool use), not from code-generation tasks specifically, and may shift for software engineering where coupling is tighter and correctness requirements stricter.

The crossover conjecture has a concrete design implication: as foundation models improve, the optimal number of parallel agents \emph{decreases}. Organizations should expect to transition from multi-agent architectures toward single-agent (or few-agent) workflows as model capability increases, redirecting coordination overhead toward richer context and more thorough verification.

\subsection{Extended Amdahl's Law with Reliability}

Classical Amdahl's Law gives speedup $S(n) = 1/(s + (1-s)/n)$ for serial fraction $s$ and $n$ parallel units. The Coordination pillar extends this with a reliability factor:

\begin{proposition}[Extended Amdahl's Law]
\label{thm:extended-amdahl}
For $n$ agents with serial fraction $s$ and per-agent defect injection rate $\delta_a$ (distinct from the context degradation function $\delta$ in Section~\ref{sec:information}), and assuming independent agent errors:
\begin{equation}
    S_{\text{eff}}(n) = \frac{1}{s + (1-s)/n} \cdot (1-\delta_a)^n
\end{equation}
The optimal agent count is $n^* \approx 1/\sqrt{\delta_a}$. Beyond $n^*$, adding agents \emph{decreases} effective throughput.
\end{proposition}

\begin{proof}[Proof sketch]
Write $\ln S_{\mathrm{eff}} = -\ln(s + (1-s)/n) + n \ln(1-\delta_a)$.
Differentiating with respect to $n$ and setting to zero:
$(1-s)/(n^2 s + n(1-s)) = -\ln(1-\delta_a) \approx \delta_a$ for small $\delta_a$.
For small $s$, this simplifies to $1/n^2 \approx \delta_a$, giving $n^* \approx 1/\sqrt{\delta_a}$.
The maximum is unique because $\ln S_{\mathrm{eff}}$ is the sum of a concave function (the Amdahl term) and a linear function with negative slope (the reliability term), hence strictly concave.
\end{proof}

At $\delta_a = 0.05$, $n^* \approx 4$ to 5 agents. At $\delta_a = 0.01$, $n^* \approx 10$. This provides a principled answer to the practitioner question ``how many agents should I run in parallel?'' and explains \citeauthor{Cursor2026}'s observation that 20 parallel agents with flat coordination collapsed to effective throughput of 2 to 3 agents.

\subsection{Universal Scalability Law Extension}

The Extended Amdahl's Law captures serialization and reliability but omits coherency overhead (the cost of keeping agents consistent with each other). Combining Gunther's Universal Scalability Law \citep{gunther2008usl} with the reliability factor gives a more complete model:

\begin{equation}
\label{eq:usl-extended}
    S_{\mathrm{eff}}(n) = \frac{n}{1 + \sigma(n-1) + \kappa \, n(n-1)} \cdot (1-\delta_a)^n
\end{equation}

where $\sigma$ is the serialization coefficient (fraction of work that must be done sequentially, analogous to $s$ in Amdahl's Law) and $\kappa$ is the coherency (crosstalk) penalty per agent pair.

The USL's $\kappa$ term is quadratic in $n$, producing a characteristic ``retrograde'' curve where throughput decreases beyond an optimal point even without the reliability factor. With the additional exponential decay from $(1-\delta_a)^n$, the optimal $n^*$ is pushed lower still, because the system faces three compounding penalties:

\begin{enumerate}[nosep]
    \item \textbf{Serialization} ($\sigma$ term): some work cannot be parallelized.
    \item \textbf{Crosstalk} ($\kappa$ term): agents must synchronize, consuming $O(n^2)$ communication overhead.
    \item \textbf{Unreliability} ($(1-\delta_a)^n$ term): each agent independently injects defects.
\end{enumerate}

These three penalties explain why empirical optimal team sizes are small: 3 to 5 agents per practitioner consensus \citep{osmani2025}, consistent with $n^* \approx 3$--$4$ from the model at typical parameter values ($\sigma \approx 0.1$, $\kappa \approx 0.02$, $\delta_a \approx 0.05$).

\subsection{Task Decomposition as Balanced Min-Cut}

The Coordination pillar formalizes task decomposition as a graph partitioning problem:

\begin{definition}[Coupling Graph]
Given a codebase with file set $\mathcal{F}$, the coupling graph $G = (\mathcal{F}, E, w)$ encodes structural imports, co-change frequency, and semantic similarity as weighted edges.
\end{definition}

Optimal decomposition for $k$ agents is a balanced $k$-way min-cut on the coupling graph, minimizing cross-partition edges (which become potential merge conflicts and coordination overhead). The Cut-Conflict Bridge lemma connects partition quality directly to expected merge conflict rate:

\begin{proposition}[Cut-Conflict Bridge]
\label{prop:cut-conflict}
For a $k$-way partition $\Pi = \{F_1, \ldots, F_k\}$ of the coupling graph, the expected merge conflict count is:
\begin{equation}
    \mathbb{E}[\text{conflicts}] = \sum_{(u,v) \in \text{cut}(\Pi)} w(u,v) \cdot p_{\text{conflict}}(u,v)
\end{equation}
where $p_{\text{conflict}}(u,v)$ is the probability that simultaneously modifying files $u$ and $v$ produces a conflict. Minimizing cut weight directly minimizes expected conflicts.
\end{proposition}

\begin{proof}[Proof sketch]
Each cross-partition edge $(u,v)$ with weight $w(u,v)$ represents coupling that may cause a conflict. The probability of a conflict arising from this edge is at most proportional to the coupling weight, by construction of the composite coupling metric (calibrated against historical conflict data). Summing over all cut edges yields the bound. The constant of proportionality depends on the calibration of $w$ against actual conflict data and is empirically in the range 0.01 to 0.1 for human-authored code; it has not been validated for agent-generated code specifically.
\end{proof}

\subsection{Five-Level Merge Conflict Taxonomy}

The Coordination pillar identifies five levels of merge conflict, ordered by detection difficulty. We present each level together with its primary detection mechanism:

\begin{table}[ht]
\centering
\caption{Five-level merge conflict taxonomy with detection methods.}
\label{tab:conflict-taxonomy}
\begin{tabular}{clll}
\toprule
\textbf{Level} & \textbf{Type} & \textbf{Detection Method} & \textbf{Detection Tool} \\
\midrule
1 & Textual & Three-way merge & git merge \\
2 & Structural & AST-level differencing & JDime, GumTree \\
3 & Dependency & Build-level analysis & Package resolver, compiler \\
4 & API & Refactoring-aware merge & IntelliMerge (88\% precision) \\
5 & Semantic & Compositional verification & SafeMerge (75\% of benchmarks) \\
\bottomrule
\end{tabular}
\end{table}

Without coordination contracts, conflict count grows as $O(k^2)$ in the number of agents (every pair can conflict). With file-ownership contracts ($F_i \cap F_j = \emptyset$), textual, structural, and API conflicts are eliminated by construction (no shared files to conflict on), reducing growth to $O(k)$ through interface-mediated dependency and semantic conflicts only. The reduction from quadratic to linear scaling is the primary formal justification for ownership-based contracts.

\subsection{Database Concurrency Control Analogy}

The Coordination pillar draws a formal analogy between multi-agent code generation and database concurrency control. Agents are transactions; files are data items; merge conflicts are write-write conflicts. The mapping is surprisingly precise:

\begin{table}[ht]
\centering
\caption{Database isolation levels mapped to agent coordination mechanisms.}
\label{tab:isolation-levels}
\begin{tabular}{llll}
\toprule
\textbf{Isolation Level} & \textbf{Agent Mechanism} & \textbf{Anomalies Prevented} & \textbf{Anomalies Permitted} \\
\midrule
Read Uncommitted & Shared workspace & None & All \\
Read Committed & Git branches, periodic pulls & Dirty reads & Non-repeatable reads, \\
 & & & phantoms, write skew \\
Snapshot & Git worktrees (frozen at & Dirty reads, non-repeatable & Write skew \\
 & dispatch) & reads, phantoms & \\
Serializable & Sequential execution & All & None \\
\bottomrule
\end{tabular}
\end{table}

The critical insight is \emph{write skew}: two agents read the same state, make individually valid changes, but the combination is invalid. Agent $A_i$ modifies file $f_a$ based on reading file $f_b$, while agent $A_j$ modifies $f_b$ based on reading $f_a$. Both commits succeed (disjoint write sets), but the combined state is inconsistent. This is the agent analog of the database write-skew anomaly, and it corresponds precisely to Level 5 (semantic) merge conflicts in the taxonomy above.

\begin{remark}
The key difference between database transactions and agent tasks is cost granularity. Database transactions are millisecond-scale; agent tasks are minute-scale. The retry cost of an aborted agent task (discarding expensive LLM inference) is orders of magnitude higher relative to task cost, shifting the calculus: even rare conflicts make optimistic approaches (detect-and-retry) less attractive compared to pessimistic approaches (prevent-by-contract) unless conflict detection happens before expensive computation.
\end{remark}

\subsection{Assume-Guarantee Composition}

Each agent $A_i$ operates under a contract $(Asm_i, Guar_i)$ following the assume-guarantee paradigm \citep{meyer1992, henzinger2002}:

\medskip
\noindent\textbf{Assumptions} $Asm_i$:
\begin{enumerate}[nosep, label=(A\arabic*)]
    \item No other agent modifies any file in $F_i$: $\forall j \neq i, \; \mathrm{WriteSet}(A_j) \cap F_i = \emptyset$.
    \item All interfaces referenced by $A_i$ but owned by other agents remain stable (signatures unchanged).
    \item The snapshot $S_i$ provided to $A_i$ at dispatch time is consistent (compiles, passes tests).
\end{enumerate}

\noindent\textbf{Guarantees} $Guar_i$:
\begin{enumerate}[nosep, label=(G\arabic*)]
    \item Agent $A_i$ modifies only files in $F_i$: $\mathrm{WriteSet}(A_i) \subseteq F_i$.
    \item For every interface $I_j$ in files owned by $A_i$, the post-modification signature matches $\mathrm{sig}_j$.
    \item The modifications, applied to $S_i$, compile and pass all tests scoped to $F_i$.
\end{enumerate}

\begin{theorem}[Contract Composition]
\label{thm:composition}
If every agent $A_i$ satisfies its contract (assuming $Asm_i$ holds, $A_i$ establishes $Guar_i$), then the composed system satisfies:
\begin{enumerate}[nosep]
    \item No write-write conflicts (from disjointness and G1).
    \item All specified interfaces are preserved (from G2).
    \item For languages with sound separate compilation (e.g., Java, Rust, Go), the merged codebase compiles (from G3 and interface preservation). For dynamically-typed languages (Python, JavaScript), the guarantee weakens to: the merged codebase compiles at all statically checkable call sites.
\end{enumerate}
\end{theorem}

\begin{proof}[Proof sketch, rely-guarantee reasoning]
\textit{Step 1 (Disjointness).}
By disjointness of $F_i$ and guarantee G1, each agent's write set is disjoint from every other agent's. This establishes the rely condition: the environment does not modify $A_i$'s files, satisfying A1.

\textit{Step 2 (Interface preservation).}
By G2, each agent preserves all interfaces in its owned files. Since A2 requires that interfaces outside $F_i$ remain stable, and G2 ensures this for every $F_j$, the circular rely-guarantee discharge holds: $Guar_j$ for $j \neq i$ collectively implies $Asm_i$.

\textit{Step 3 (Compilation).}
Each agent's modifications compile in isolation (G3) with all consumed interfaces preserved (A2, now discharged). Since the merged codebase preserves all interfaces, every call site remains type-consistent. By compositionality of type checking, the merged result type-checks.
\end{proof}

\begin{remark}
This proof is modular: adding agent $A_{k+1}$ with a fresh disjoint $F_{k+1}$ requires only verifying $A_{k+1}$'s contract, not re-verifying existing agents. This modularity is the key advantage of assume-guarantee reasoning over monolithic verification.
\end{remark}

\subsection{Conway's Law as Graph Containment}

The Coordination pillar formalizes Conway's Law as a graph-theoretic theorem. Define the \emph{communication graph} $G_T = (\mathcal{A}, E_T)$ where an edge $(A_i, A_j) \in E_T$ indicates that agents $A_i$ and $A_j$ can exchange messages under topology $T$. Define the \emph{required communication graph} $G_R = (\mathcal{A}, E_R)$ where $(A_i, A_j) \in E_R$ indicates that agents $A_i$ and $A_j$ manage modules with inter-module dependencies.

\begin{theorem}[Conway's Law, Graph-Theoretic Form]
\label{thm:conway}
If $G_R \not\subseteq G_T$ (the required communication graph is not a subgraph of the topology communication graph), then for every edge $(A_i, A_j) \in E_R \setminus E_T$, the agents managing the dependent modules cannot coordinate, producing a semantic conflict with probability proportional to the coupling strength of the underlying module dependency.
\end{theorem}

The condition for low-overhead coordination is $G_R \subseteq G_T$: every strong coupling between files must be matched by a communication channel between the agents responsible for those files. When this condition fails, the ``topology-module mismatch cost'' is:

\begin{equation}
    \mathrm{Mismatch}(\alpha, T, G_M) = \sum_{(A_i, A_j) \in E_R \setminus E_T} \sum_{\substack{(M_a, M_b) \in E_M \\ \alpha(M_a)=A_i, \alpha(M_b)=A_j}} w(M_a, M_b)
\end{equation}

where $\alpha: \mathcal{M} \to \mathcal{A}$ is the module-to-agent assignment. Mismatch $= 0$ when $E_R \subseteq E_T$ (Conway-aligned).

\subsection{Inverse Conway for Agents}

Conway's Law describes how organizational structure constrains system structure. The \emph{Inverse Conway Maneuver} asks: given a desired system structure, how should we organize the agents?

\begin{theorem}[Inverse Conway for Agents]
\label{thm:inverse-conway}
The assignment $\alpha^*$ minimizing mismatch for a given topology $T$ is the solution to a constrained balanced min-cut: partition the module graph into $k$ parts minimizing uncovered cross-partition edge weight.
\end{theorem}

For the independent topology ($E_T = \emptyset$), mismatch equals total cross-partition coupling, reducing to standard balanced min-cut. For mesh topology ($E_T$ is complete), mismatch is zero for any assignment, but communication overhead is $O(k^2)$. For isolated-then-merge with contracts at coverage $\gamma$, the effective mismatch is $(1-\gamma)$ times the full mismatch, since contracts act as static communication channels conveying interface information without runtime message passing.

This connects task decomposition (Section 5.6) to topology selection: the optimal decomposition depends on the topology, and the optimal topology depends on the decomposition. In practice, this joint optimization is solved iteratively: choose a topology, compute the best decomposition for that topology, evaluate QAS, and compare across topologies.

\subsection{Quality-Adjusted Speedup}

\begin{definition}[Quality-Adjusted Speedup]
\label{def:qas}
\begin{equation}
    \text{QAS}(k) = \frac{T_1 \cdot Q_1}{T_k \cdot Q_k}
\end{equation}
where $T_k$ is wall-clock time with $k$ agents and $Q_k$ is output quality (e.g., fraction of changes that would pass expert review). QAS $< 1$ means parallelism is net-negative: the quality degradation outweighs the speed improvement.
\end{definition}

Calibration against DORA 2025 data suggests that naive AI-assisted parallelism frequently yields QAS $< 1$: individual task completion up 21\%, but review time up 91\%, PR size up 154\%, and organizational delivery metrics flat despite individual gains. (The ``bug rates up 9\%'' figure appearing in secondary analyses could not be independently verified from the primary DORA report.)

More precisely, for any topology $t$ and agent count $k > 1$, if the error amplification factor $A_e(t,k)$ satisfies $A_e(t, k) \cdot d(1) \geq 1 - (1 - d(1))/S(k)$ (where $d(1)$ is the single-agent defect rate and $S(k)$ is the raw speedup), then QAS$(k) < 1$. This condition is easily satisfied when $d(1)$ is moderate (0.15 to 0.30), $A_e$ is large (4 to 17), and $S(k)$ is modest (2 to 3). The regime QAS $< 1$ is not a corner case; it is the default outcome for poorly coordinated multi-agent systems.


% ====================================================================
% 6. THE TEMPORAL PILLAR
% ====================================================================
\section{The Temporal Pillar}
\label{sec:temporal}

The Temporal pillar addresses the time dimension of harness architecture: how fast should agents iterate, when should they speculate, and how should they manage information freshness?

\subsection{Context Fidelity Decay: The Channel Model}

We model the agent's context as a noisy channel between the true codebase state $\mathcal{X}$ and the agent's representation $\mathcal{Y}$. At the moment of a fresh context read ($t = t_0$), the channel is noiseless: $I(\mathcal{X}; \mathcal{Y}) = I_{\mathrm{fresh}}$. As time passes, the codebase evolves while the agent's snapshot remains fixed, introducing noise.

\begin{definition}[Context Fidelity]
\label{def:fidelity}
The context fidelity $\Phi(t)$ at time $t$ for a context snapshot taken at $t_0$ is:
\begin{equation}
    \Phi(t) = \frac{I_{\mathrm{stale}}(t)}{I_{\mathrm{fresh}}} = e^{-\Lambda(t - t_0)}
\end{equation}
where $\Lambda = \lambda \alpha$ is the effective decay rate, $\lambda$ is the codebase change rate (changes per unit time), and $\alpha$ is the average information disruption per change.
\end{definition}

\begin{theorem}[Exponential Staleness Decay]
\label{thm:staleness}
Assume: (A1) codebase changes follow a Poisson process with rate $\lambda$; (A2) each change independently invalidates a fraction $\alpha$ of the agent's cached context; (A3) $\alpha$ is constant across change events. Then:
\begin{equation}
    I_{\mathrm{stale}}(t) = I_{\mathrm{fresh}} \cdot e^{-\lambda\alpha(t - t_0)}
\end{equation}
\end{theorem}

\begin{proof}
The number of changes in $[t_0, t]$ is $N \sim \mathrm{Poisson}(\lambda\Delta t)$ where $\Delta t = t - t_0$. After $k$ changes, the surviving information fraction is $(1-\alpha)^k$. Taking the expectation via the Poisson probability generating function:
\begin{equation}
    \mathbb{E}[(1-\alpha)^N] = e^{\lambda\Delta t((1-\alpha) - 1)} = e^{-\lambda\alpha\Delta t}
\end{equation}
\end{proof}

The context fidelity has half-life $t_{1/2} = \ln 2 / \Lambda$.

The exponential decay model is calibrated against the 17.1\% ``Fresh Eyes'' advantage reported by \citet{suzgun2024metaprompting}: agents with fresh context outperform agents with stale context by this margin. Treating this as an upper bound on fidelity loss (the advantage likely conflates freshness with task decomposition effects):
\begin{equation}
    1 - \Phi_{\mathrm{inc}} \leq 0.171 \implies \Phi_{\mathrm{inc}} \geq 0.829
\end{equation}

With a typical cycle of $\bar{t} \approx 50$ minutes $\approx 0.83$ hours, this gives an upper bound on the effective decay rate of $\Lambda \leq 0.226$ per hour (context half-life $\geq 3$ hours).

\begin{remark}
Assumptions A2 and A3 are strong. In practice, changes are correlated (a refactor touches many files; a bugfix touches one), and the disruption fraction $\alpha$ varies across changes. The single-rate model should be understood as a first-order approximation. A multi-exponential extension partitions context into strata with different decay rates: $I_{\mathrm{stale}}(t) = \sum_{k=1}^K I_k \cdot e^{-\lambda\alpha_k(t - t_0)}$, capturing the empirical observation that architectural knowledge ($\alpha_k$ small) persists longer than implementation details ($\alpha_k$ large).
\end{remark}

In multi-agent settings, $\Lambda$ increases with the number of concurrent agents (more agents means faster codebase evolution), creating a feedback loop between parallelism and context staleness. This connects directly to the Coordination pillar's optimal agent count: adding agents increases $\lambda$ (the codebase change rate), which increases $\Lambda$, which decreases $\Phi$, which decreases per-agent quality, which further reduces the effective speedup.

\subsection{Verified Iterations per Hour}

\begin{definition}[Verified Iterations per Hour]
\label{def:vih}
\begin{equation}
    \text{VIH} = \lambda_{\text{raw}} \cdot p_{\text{verify}}
\end{equation}
where $\lambda_{\text{raw}}$ is the raw iteration rate (iterations per hour) and $p_{\text{verify}}$ is the probability that an iteration passes all verification gates.\footnote{VIH has not, to our knowledge, been measured in any production harness. It is proposed here as a metric, not reported as an empirical finding.}
\end{definition}

VIH captures the fundamental speed-quality tradeoff. Skipping verification increases $\lambda_{\text{raw}}$ but decreases $p_{\text{verify}}$. The Temporal pillar proves that VIH is maximized at the unit-elasticity point:

\begin{proposition}[VIH Optimality]
\label{prop:vih-optimal}
VIH is maximized when:
\begin{equation}
    \frac{\partial \ln \lambda_{\text{raw}}}{\partial v} = -\frac{\partial \ln p_{\text{verify}}}{\partial v}
\end{equation}
where $v$ is the verification investment. At this point, the marginal throughput gain from reducing verification exactly equals the marginal quality loss.
\end{proposition}

\begin{proof}[Proof sketch]
Differentiate $\text{VIH}(\lambda) = \lambda \cdot p_{\mathrm{verify}}(\lambda)$ and set to zero: $p_{\mathrm{verify}}(\lambda) + \lambda \cdot p_{\mathrm{verify}}'(\lambda) = 0$. Rearranging: $p_{\mathrm{verify}}'(\lambda^*)/p_{\mathrm{verify}}(\lambda^*) = -1/\lambda^*$, which is the unit-elasticity condition. The optimum occurs where a 1\% increase in raw speed causes exactly a 1\% decrease in verification rate.
\end{proof}

Under an exponential degradation model $p_{\mathrm{verify}}(\lambda) = p_0 e^{-\beta(\lambda - \lambda_0)}$ (chosen for analytical convenience; the true functional form is unknown), the optimal rate is $\lambda^* = 1/\beta$.

\subsection{Amdahl's Law Bound on VIH}

If verification is a serial fraction $f$ of the cycle (it runs after execution and cannot be overlapped), then even with infinite parallelization of all other stages:

\begin{equation}
    \text{VIH}^{\max} = \frac{p_{\text{verify}}}{f \cdot \tau_{\text{cycle}}}
\end{equation}

This is a direct application of Amdahl's Law \citep{amdahl1967validity}: no amount of speedup in non-verification stages can overcome the serial verification bottleneck.

From representative harness data, $f \approx 0.20$ (verification takes roughly 10 of 50 minutes). With $p_{\text{verify}} \approx 0.67$ (using Devin's PR merge rate as a proxy; \citealt{cognition2025devin}), $\text{VIH}^{\max} \approx 4.0$ verified iterations per hour, regardless of non-verification speedups.

\begin{remark}
The Amdahl bound implies that the highest-leverage temporal optimization is reducing verification time while maintaining verification quality, not speeding up any other pipeline stage. Tiered verification (fast static checks first, expensive dynamic checks only when static checks pass) is the natural strategy.
\end{remark}

\subsection{The Speculation Ratio Test}

Speculative execution (beginning the next pipeline stage before the current stage's output is verified) has positive expected value when:

\begin{theorem}[Speculation Ratio Test]
\label{thm:speculation}
Speculate if and only if:
\begin{equation}
    \frac{p}{q} > \frac{R}{B}
\end{equation}
where $p$ is the probability that the current stage's output is correct, $q = 1 - p$ is the failure probability, $R$ is the cost of rollback on failure, and $B$ is the benefit of early completion on success.
\end{theorem}

\begin{proof}[Proof sketch]
Expected speculative time: $\mathbb{E}[T_{\mathrm{spec}}] = p \cdot \mathbb{E}[\max(t_i, t_j)] + q(R + \mathbb{E}[t_j])$. Sequential baseline: $\mathbb{E}[T_{\mathrm{seq}}] = \mathbb{E}[t_i] + \mathbb{E}[t_j]$. The saving under success is $B = \mathbb{E}[t_i] + \mathbb{E}[t_j] - \mathbb{E}[\max(t_i, t_j)]$. Speculation is profitable when $p \cdot B > q \cdot R$, which rearranges to the stated condition.
\end{proof}

For typical agent pipelines with $p \approx 0.85$ to $0.95$ and $R/B \approx 0.3$ to $0.5$, speculation is usually beneficial for the first stage ahead but not for deeper speculation (where rollback costs compound).

\subsection{Speculation Depth Limit}

For multi-stage pipelines, speculation can be applied at multiple depths (speculating two, three, or more stages ahead). However, the probability of a full chain succeeding decays exponentially, while rollback costs grow.

\begin{theorem}[Speculation Depth Limit]
\label{thm:depth-limit}
For a homogeneous chain with per-stage success probability $p$, benefit $B_0$, and rollback cost $R_0$, the maximum useful speculation depth is:
\begin{equation}
    k^* < \frac{\ln(1 + B_0/R_0)}{\ln(1/p)}
\end{equation}
\end{theorem}

\begin{proof}[Proof sketch]
At depth $k$, the probability of needing rollback is $1 - p^k$, and the accumulated rollback cost is at most $k \cdot R_0$. The benefit is at most $k \cdot B_0 \cdot p^k$ (all $k$ stages must succeed). Setting the net expected value to zero: $p^k \cdot k B_0 = (1 - p^k) \cdot k R_0$. Solving for $p^k$: $p^k > R_0/(R_0 + B_0)$. Taking logarithms yields the stated bound.
\end{proof}

For $p = 0.85$ and $B_0/R_0 = 2$: $k^* \leq 6$ stages. For $p = 0.70$ and $B_0/R_0 = 1$: $k^* \leq 1$. The exponential decay of $p^k$ places a natural horizon on useful speculation, consistent with the practical finding that deep speculation chains in agent pipelines are rarely profitable.

\begin{remark}[The Spectre Analogy]
Speculative execution in CPUs was considered a pure optimization for two decades before Spectre and Meltdown revealed that speculative state changes created security vulnerabilities. In agent pipelines, the analogous risks include: state pollution from discarded speculative work, combinatorial branching complexity, debugging difficulty when causal chains include hypothetical branches, and cognitive overhead for humans monitoring speculative pipelines. These hidden costs suggest that speculation should be applied conservatively, typically at depth 1.
\end{remark}

\subsection{Three Caching Layers}

Agent harnesses interact with three distinct caching layers, each governed by different dynamics and subject to different staleness tradeoffs:

\begin{enumerate}
    \item \textbf{KV-Cache} (attention state reuse, LRU-governed): Provider-level caching of computed key-value pairs for shared prompt prefixes. Eviction follows LRU (Least Recently Used) policy at the provider level. Claude Code achieves 92\% hit rate and 81\% cost reduction (vendor-reported; \citealt{anthropic2026agentic}). Hit rate depends on prompt structure: append-only prompts with static prefixes maximize hits.

    \item \textbf{Prompt/Semantic Caching} (token-level, semantic similarity): Server-side caching of prompt computation for repeated or semantically similar prefixes. \citet{liang2026cache} report 41--80\% cost reduction and 13--31\% time-to-first-token improvement across providers. Semantic caching extends exact-match caching by allowing fuzzy retrieval, trading match precision for hit rate.

    \item \textbf{Plan Template Caching} (reasoning-level, persistent): Reuse of structured plan templates across similar tasks. \citet{zhang2025apc} achieve 50\% cost and 27\% latency reduction while maintaining 96.6\% of optimal performance. Templates persist across sessions and are the most durable cache layer, but also the most vulnerable to staleness (a plan template designed for one codebase version may be subtly wrong for the next).
\end{enumerate}

The Denning working set model \citep{denning1968working} applies to all three layers: cache effectiveness depends on temporal locality in the request stream. The cache TTL must match the agent's task locality window to avoid thrashing.

\subsection{Cache Competitive Analysis}

The Temporal pillar formulates context caching as an inventory problem, where cache entries are ``inventory'' that becomes stale over time:

\begin{theorem}[Optimal Refresh Interval]
\label{thm:cache-eoq}
The optimal cache refresh interval balancing refresh cost $R$ against staleness cost is:
\begin{equation}
    T^* = \sqrt{\frac{2R}{\lambda \alpha \cdot E_{\text{stale}}}}
\end{equation}
where $\lambda$ is the access rate, $\alpha$ is the staleness growth rate, and $E_{\text{stale}}$ is the expected cost per stale access. This is an EOQ (Economic Order Quantity) analog for information freshness.
\end{theorem}

\begin{proof}
The amortized cost per unit time is $C(T) = R/T + \frac{1}{2}\lambda\alpha T \cdot E_{\mathrm{stale}}$ (the staleness integral over $[0, T]$ under linear growth averages to $\frac{1}{2}\lambda\alpha T$). Setting $C'(T) = 0$: $-R/T^2 + \frac{1}{2}\lambda\alpha E_{\mathrm{stale}} = 0$, yielding the stated formula.
\end{proof}

The competitive analysis from the online algorithms literature provides worst-case guarantees for cache refresh policies:

\begin{itemize}[nosep]
    \item The break-even deterministic policy (refresh after $\lceil R/r \rceil$ incremental steps, where $r$ is the incremental update cost) is 2-competitive against the offline optimal \citep{sleator1985amortized}. That is, its total cost is at most twice the cost of the optimal policy that knows the future request sequence.
    \item The optimal randomized policy achieves $e/(e-1) \approx 1.58$-competitive \citep{fiat1991competitive}.
\end{itemize}

The cache introduces a fundamental tension: append-only context maximizes cache hits but accumulates noise, while context reset loses the cache but gains fresh reasoning quality (the 17.1\% advantage). The optimal cycle length is partially determined by when cache savings from continuation are outweighed by quality loss from context accumulation.

\subsection{Bayesian Mode Selection with Bainbridge's Ironies}

The Temporal pillar provides a Bayesian decision framework for selecting between execution modes:

\begin{itemize}[nosep]
    \item \textbf{Fast mode}: minimal verification, high $\lambda_{\text{raw}}$, low $p_{\text{verify}}$. Optimal for low-risk, well-understood tasks.
    \item \textbf{Governed mode}: full verification, low $\lambda_{\text{raw}}$, high $p_{\text{verify}}$. Optimal for high-risk, novel tasks.
\end{itemize}

\begin{theorem}[Bayesian Mode Selection]
\label{thm:bayesian-mode}
For the binary simplification (Fast vs.\ Governed), the optimal decision rule is:
\begin{equation}
    \text{Choose Governed iff } P(\text{Governed needed} \mid \mathbf{x}) > \frac{c_{\mathrm{GF}}}{c_{\mathrm{FG}} + c_{\mathrm{GF}}}
\end{equation}
where $c_{\mathrm{FG}}$ is the cost of choosing Fast when Governed is needed (undetected quality failure, typically 10--100$\times$ the task time in rework) and $c_{\mathrm{GF}}$ is the cost of choosing Governed when Fast suffices (wasted time, typically 2--5$\times$ the task time).
\end{theorem}

Since $c_{\mathrm{FG}} \gg c_{\mathrm{GF}}$, the threshold is small (typically 0.05--0.10), meaning the harness should choose Governed even when the probability of needing it is quite low. This is the standard asymmetric-cost Bayesian decision rule applied to the mode selection problem.

However, \citet{bainbridge1983ironies} identified a double bind in automated decision-making that complicates this framework: if mode selection rules can be fully specified, a computer can apply them faster than a human monitor can verify; if mode selection requires judgment, automation will fail at boundary cases while the human monitor, deprived of practice, performs even worse. The implication for harness design is threefold:

\begin{enumerate}[nosep]
    \item The mode selection classifier must be \emph{conservative} (biased toward Governed), because the cost of misclassification is asymmetric.
    \item The classifier should operate as a \emph{proposal} that the human confirms or overrides, rather than as a fully autonomous decision.
    \item The feature engineering problem is recursive: deciding how much analysis to do for mode classification requires analysis, consuming the very budget that mode selection is meant to allocate efficiently.
\end{enumerate}

\subsection{The Speed-Quality Pareto Frontier}

In the space of $(\lambda_{\text{raw}}, p_{\text{verify}})$, achievable operating points form a Pareto frontier. The frontier has an empirically motivated three-regime structure:

\begin{enumerate}[nosep]
    \item \textbf{Concave at low speed}: redundant checks provide diminishing quality returns; easy to gain speed cheaply by removing the lowest-value verification.
    \item \textbf{Convex at moderate speed}: each further speed gain requires skipping progressively more important verification, with accelerating quality loss.
    \item \textbf{Concave again at high speed}: statistical filtering via best-of-$N$ stabilizes quality (running $5\times$ more cheap iterations with majority voting can yield higher effective quality than fewer careful iterations, per SWE-bench data showing Pass@5 substantially exceeding Pass@1; \citealt{sweeffi2025}).
\end{enumerate}

\subsection{Pareto Frontier Mixing Strategies and Convexification}

\begin{theorem}[Convexification by Mixing]
\label{thm:convexification}
In concave frontier regions, the mixed strategy using configuration $c_1$ with probability $\alpha$ and $c_2$ with probability $(1-\alpha)$ achieves:
\begin{equation}
    (S_{\mathrm{mix}}, Q_{\mathrm{mix}}) = \alpha(S_1, Q_1) + (1-\alpha)(S_2, Q_2)
\end{equation}
which Pareto-dominates every pure strategy in that region. The set of achievable operating points is the \emph{convex hull} of the pure-strategy frontier, not the frontier itself.
\end{theorem}

This is the direct analogue of the classical result that mixed strategies expand the achievable set to the convex hull of pure-strategy outcomes. The practical implication is that a harness alternating between ``fast mode for 70\% of tasks'' and ``governed mode for 30\% of tasks'' can achieve a speed-quality combination unattainable by any single fixed configuration.

\begin{theorem}[VIH Tangency Condition]
\label{thm:vih-tangency}
VIH $= S \cdot Q$ is maximized on the frontier where the iso-VIH hyperbola $Q = \text{VIH}/S$ is tangent to the frontier curve $Q = f(S)$. At this point, the elasticity of quality with respect to speed is exactly $-1$:
\begin{equation}
    \varepsilon_{Q,S} = \frac{dQ}{dS} \cdot \frac{S}{Q} = -1
\end{equation}
Below this point, speed gains improve VIH; above it, speed gains degrade VIH.
\end{theorem}

\begin{proof}
Maximize $\text{VIH}(S) = S \cdot f(S)$. Setting $\text{VIH}'(S) = f(S) + S f'(S) = 0$ gives $f'(S^*) = -f(S^*)/S^* = -Q^*/S^*$, which is the unit-elasticity condition.
\end{proof}

The convexification theorem implies that the VIH-optimal operating point on the convexified frontier may correspond to a mixing strategy rather than a pure configuration. Specifically, if the VIH-maximizing tangent point lies in a concave region of the original frontier, the optimal policy is to randomize between the two pure strategies at the endpoints of that concave segment. This connects execution mode selection (choosing Fast vs. Governed per task) to Pareto frontier optimization: the Bayesian mode selector is implementing the mixing strategy that achieves the VIH-optimal point on the convexified frontier.
% ====================================================================
% 7. THE QUALITY PILLAR (EXPANDED)
% ====================================================================
\section{The Quality Pillar}
\label{sec:quality}

The Quality pillar addresses the specific defects introduced by AI code generation and the formal structure of defense against them. The central insight is that AI-generated code produces a qualitatively new class of defect: code that is superficially competent, individually defensible, and cumulatively corrosive to codebase health. We call these defects \emph{slop}, following industry usage.

\subsection{Slop Taxonomy: 18 Types in 3 Decidability Classes}

The Quality pillar classifies 18 AI code defect types (``slop types'') into three decidability classes, grounded in Rice's theorem. Each slop type is formalized as a tuple $\tau_j = (\sigma_j, \phi_j, \delta_j, C_j)$ where $\sigma_j$ is a syntactic signature on ASTs, $\phi_j$ is the semantic property determining genuine defect status given full program context, $\delta_j \in \{\mathrm{D}, \mathrm{SD}, \mathrm{U}\}$ is the detection decidability class, and $C_j$ is the cost distribution over deployment contexts.

\begin{table}[htbp]
\centering
\caption{The 18 slop types partitioned by detection decidability.}
\label{tab:slop-partition}
\small
\begin{tabularx}{\textwidth}{l c X}
\toprule
\textbf{Class} & \textbf{Count} & \textbf{Types} \\
\midrule
D (Decidable) & 7 & Hallucinated imports, placeholder/stub code, duplicate logic, dead code, lint suppressions, cross-language contamination, outdated APIs \\
\addlinespace
SD (Semi-decidable) & 7 & Security vulns (functional-looking), happy-path error handling, data model mismatches, over-engineering, architectural erosion, silent scope expansion, god functions \\
\addlinespace
U (Undecidable) & 4 & Meaningless tests, boilerplate inflation, redundant comments, naming/convention drift \\
\bottomrule
\end{tabularx}
\end{table}

The three classes are formally distinct:

\begin{enumerate}
    \item \textbf{Decidable (D):} There exists a total computable function $f_j$ such that $f_j(t) = 1$ iff $t$ contains an instance of $\tau_j$. The syntactic signature $\sigma_j$ suffices; the semantic property $\phi_j$ does not depend on unobservable context. Hallucinated imports are detected by lockfile resolution ($O(1)$ table lookup); placeholder code by AST pattern matching for \texttt{pass}, \texttt{...}, \texttt{NotImplementedError}; duplicate logic by tree-sitter-based clone detection (suffix tree, $O(n \log n)$); dead code by call graph reachability analysis.

    \item \textbf{Semi-decidable (SD):} The semantic property $\phi_j$ depends on partially formalizable context. An oracle (LLM or human) can approximate $\phi_j$ with probability $p \in (0.5, 1)$, but no computable function achieves $p = 1$. Security vulnerabilities require knowing which inputs are adversarial; happy-path handling requires knowing which error paths matter; over-engineering requires judging abstraction necessity. Expert inter-rater agreement is $\kappa \approx 0.6$ (substantial, above chance but below certainty).

    \item \textbf{Undecidable (U):} No known oracle achieves reliable detection. For all known procedures $f$, inter-rater reliability $\kappa < 0.4$ (below ``fair agreement''). ``Meaningless test'' depends on test intent versus structure; ``boilerplate'' thresholds vary by team norms; ``redundant comment'' depends on reader expertise; ``naming drift'' requires inferring implicit conventions.
\end{enumerate}

This classification has a direct design implication: it is impossible to build a fully automated quality gate that catches all slop types. The four-layer verification architecture (shared with the Reliability pillar) is a necessary response to this computational limitation.

\subsection{Three Generative Factors}

The 18 slop types arise from three approximately orthogonal generative factors operating at different stages of the generation process:

\begin{definition}[Context Blindness (F1)]
The agent's effective context $\hat{\mathcal{K}}(t) \subset \mathcal{K}(t)$, where $\mathcal{K}(t)$ is the codebase knowledge required for correct generation at point $t$. The probability of slop type $\tau_j$ increases monotonically with the context deficit $|\mathcal{K} \setminus \hat{\mathcal{K}}|$. Types loading on F1: duplicate logic, naming drift, architectural erosion, boilerplate inflation.
\end{definition}

\begin{definition}[Completion Bias (F2)]
The agent's output distribution $\pi$ assigns higher probability to syntactically complete outputs than to semantically minimal correct outputs: $\mathbb{E}_\pi[\mathrm{length}(o)] > \mathbb{E}_{\pi^*}[\mathrm{length}(o)]$ where $\pi^*$ is optimal. Types loading on F2: placeholder code, happy-path error handling, meaningless tests, redundant comments, over-engineering.
\end{definition}

\begin{definition}[Training Distribution Leakage (F3)]
Patterns memorized during training are applied to contexts where they do not hold. Probability increases with divergence between training distribution $\mathcal{D}$ and execution environment $\mathcal{E}$. Types loading on F3: hallucinated imports, cross-language contamination, outdated APIs, security vulnerabilities.
\end{definition}

If the three factors are approximately orthogonal ($I(F_i; F_k)$ small relative to $H(F_i)$ for $i \neq k$), then interventions targeting different factors compose multiplicatively. Better retrieval (addressing F1) does not reduce completion bias (F2), and vice versa. The quality stack must address all three independently.

\subsection{Layered Defense Optimization}

Given 18 slop types $J = \{1, \ldots, 18\}$ and four defense layers $L = \{1, 2, 3, 4\}$ (structural gates, deterministic analysis, LLM-as-Judge, human review), the problem is to select $\Lambda_j \subseteq L$ for each type $j$ minimizing total cost:
\begin{equation}
C_{\mathrm{total}} = \sum_{\ell \in L} c_\ell \cdot \mathbf{1}[\Lambda^\ell \neq \emptyset] + \sum_{j \in J} p_j \cdot D_j \cdot P_{\mathrm{esc}}(j, \Lambda_j)
\end{equation}
where $P_{\mathrm{esc}}(j, \Lambda_j) = P(\bigcap_{\ell \in \Lambda_j} \{Z_{\ell,j} = 0\})$ is the escape probability under correlation, and $\Lambda^\ell = \{j : \ell \in \Lambda_j\}$.

\begin{theorem}[Layered Defense NP-Hardness]
\label{thm:defense-nphard}
The problem of assigning slop types to defense layers to minimize total cost subject to a minimum detection rate constraint is NP-hard, via reduction from Weighted Set Cover.
\end{theorem}

\begin{proof}[Proof sketch]
Set $D_j = M$ for large $M$ and $p_j = 1$ for all $j$. Then minimizing the objective forces all types to be covered (the penalty $M \cdot P_{\mathrm{esc}} \gg c_\ell$ for any uncovered type), reducing to minimum-weight subcollection selection, which is Weighted Set Cover. The inapproximability threshold is $\ln |J|$ unless P $=$ NP.
\end{proof}

\begin{proposition}[Greedy Approximation]
\label{prop:greedy-defense}
The greedy algorithm (assign each slop type to its highest-efficiency layer first, then fill coverage gaps) achieves a $(1-1/e)$-approximation, since the objective function is submodular.
\end{proposition}

\begin{proof}[Proof sketch]
Under independence, the marginal escaped-defect-cost reduction from adding a layer is monotone submodular (adding a layer never increases escape probability, and marginal reduction diminishes as more layers are selected). The result follows from \citet{nemhauser1978}.
\end{proof}

\begin{remark}
For our problem ($|L| = 4$, $|J| = 18$), exact enumeration over $2^4 = 16$ subsets per type is computationally trivial. The greedy bound matters for scaled versions with dozens of specialized tools.
\end{remark}

\subsection{Connection to Littlewood-Strigini and the Diversity Gain}

\citet{littlewood1993, littlewood2000} proved that independently developed software versions do not fail independently, because some inputs are inherently difficult. Their result adapts directly to LLM judge-producer pairs: if producer and judge failure probabilities both increase with difficulty, then $P(\text{defect produced and missed}) > P(\text{produced}) \cdot P(\text{missed})$.

\begin{proposition}[Correlated Judge-Producer Errors]
\label{prop:fkg}
Under the FKG inequality, if the producer's failure event and the judge's failure event are positively correlated (both more likely for similar inputs), then:
\begin{equation}
    P(\text{both miss defect}) \geq P(\text{producer misses}) \cdot P(\text{judge misses})
\end{equation}
The independence assumption \emph{underestimates} the probability of undetected defects.
\end{proposition}

This result formalizes why LLM-as-Judge systems using the same model family as the code producer underperform expectations. The \emph{diversity gain} from cross-model judging (using judge $M_J^B$ from a different vendor than producer $M_P^A$) versus same-model judging ($M_J^A$) is:
\begin{equation}
\Delta_{\mathrm{div}} = P_{\mathrm{esc}}^{A \text{ judges } A} - P_{\mathrm{esc}}^{B \text{ judges } A}
\end{equation}
This gain has two components: the detection rate difference ($d_J^B$ versus $d_J^A$) and the correlation reduction ($\rho_{AB} < \rho_{AA}$). Even at equal detection rates, diversity gain is positive whenever $\rho_{AB} < \rho_{AA}$.

Empirically, diverse techniques can achieve effectiveness \emph{greater} than the independence assumption predicts (positive diversity gain), while same-technique repetitions underperform it \citep{littlewood2000}.

\subsection{Optimal Layer Ordering}

\begin{theorem}[ROI Ordering]
\label{thm:roi-ordering}
Under independence, when layers form a cascade (resolved items skip subsequent layers), cost-optimal ordering sorts by decreasing $r_\ell / c_\ell$, where $r_\ell$ is the resolution rate.
\end{theorem}

\begin{proof}[Proof sketch]
Pairwise exchange argument (Smith's rule): swapping adjacent layers $\ell_a, \ell_b$ shows the ordering $\ell_a$ before $\ell_b$ is cheaper iff $r_a / c_a > r_b / c_b$.
\end{proof}

Under positive correlation, ROI ordering can be suboptimal. The marginal value of a layer depends more on its independence from existing layers than on its absolute detection rate.

\subsection{The Turing Enforcement Principle}

Rice's theorem establishes that for any non-trivial semantic property $P$ of programs (depending on input-output behavior, not syntactic form), the set $\{e : \phi_e \in P\}$ is undecidable. This motivates a fundamental enforcement gap:

\begin{theorem}[Enforcement Gap]
\label{thm:enforcement-gap}
For any syntactic property $P_S$, there exists a deterministic enforcer with $P(\text{correct enforcement}) = 1$. For any semantic property $P_U$, every enforcer (including any LLM) satisfies $P(\text{correct enforcement}) < 1$. The gap cannot be closed by prompt engineering alone.
\end{theorem}

\begin{proof}[Proof sketch]
Part 1: syntactic properties are decided by finite automata on the AST, achieving zero error probability. Part 2: an LLM is a fixed-weight transformer with finite training data. Rice's theorem establishes that no algorithm decides $P_U$ for all programs; the practical consequence is that achieving $P = 1$ would require correctly classifying programs from an unbounded space, including those arbitrarily far from the training distribution.
\end{proof}

The AgentSpec result \citep{wang2026agentspec} provides empirical calibration: 87.26\% enforcement for code-based rules versus 77\% for prompt-based rules. The 10-point gap is the measured enforcement gap.

\subsection{Detection Ceiling via Data Processing Inequality}

\begin{theorem}[Detection Ceiling]
\label{thm:detection-ceiling}
For any defense layer $\ell$ with access to observable features $X$ (code text, AST, test results), the detection rate for defect type $j$ satisfies:
\begin{equation}
    d_{\ell,j} \leq \frac{I(X; D_j)}{H(D_j)}
\end{equation}
where $D_j$ is the binary indicator of defect presence and $I(X; D_j)$ is the mutual information between observable features and defect status. This follows from the Data Processing Inequality.
\end{theorem}

The Detection Ceiling makes precise the intuition that some defects cannot be reliably detected from code alone. When the mutual information $I(X; D_j)$ is low (because the defect is about intent, not syntax), no amount of analytical sophistication can raise detection above the ceiling. For example, silent scope expansion requires the task description as context; Layer 1 (structural gates) does not observe the task description, so $d_{1,13} \approx 0$.

\subsection{Codebase Entropy and the Accretion Category}

\begin{definition}[Codebase Entropy]
\label{def:codebase-entropy}
For a codebase $\mathcal{C}$ of $n$ files, let $\mu(f_i)$ be the minimum set of files a developer must understand to correctly modify $f_i$. The codebase entropy is:
\begin{equation}
H(\mathcal{C}) = \frac{1}{n} \sum_{i=1}^{n} \log_2 |\mu(f_i)|
\end{equation}
A perfectly modular codebase has $H(\mathcal{C}) = 0$; a fully coupled one has $H(\mathcal{C}) = \log_2 n$.
\end{definition}

\begin{definition}[Accretion Rate]
The accretion rate $\alpha(t) = (H(\mathcal{C}_{t+1}) - H(\mathcal{C}_t)) / |\Delta_t|$ measures entropy increase per unit of code change.
\end{definition}

\begin{theorem}[Compound Degradation]
\label{thm:compound-degradation}
A sequence of $n$ individually CI-passing changes, each adding $\Delta H$ bits of entropy (unnecessary complexity, duplicated logic, expanded scope), produces total entropy growth:
\begin{equation}
    H_{\text{total}} \geq n \cdot \Delta H + \binom{n}{2} \cdot \Delta H_{\text{cross}}
\end{equation}
where $\Delta H_{\text{cross}} \geq 0$ accounts for cross-change interactions. The growth is superlinear when $\Delta H_{\text{cross}} > 0$.
\end{theorem}

\begin{proof}[Proof sketch]
Each change adds code without refactoring, increasing both $n$ and coupling degree. The sum telescopes, and the $\log_2 |\mu(f_i)|$ term grows as $\beta \log_2 n$ under the assumption $|\mu(f_i)| \sim n^\beta$ with $0 < \beta < 1$. CI checks are \emph{pointwise} (verifying the current change) but entropy is \emph{global} (depending on cumulative coupling). Pointwise correctness does not imply bounded global entropy.
\end{proof}

This theorem explains how codebases degrade under AI-generated code even when every individual change passes CI. The degradation is not visible at the per-change level; it accumulates through interactions between changes, manifesting as increasing complexity, decreasing coherence, and rising bug rates.

\subsection{The DRE Cascade Theorem}

\begin{theorem}[DRE Cascade]
\label{thm:dre-cascade}
For $L$ independent layers, cumulative Defect Removal Efficiency for type $j$ is:
\begin{equation}
D_j = 1 - \prod_{\ell=1}^{L} (1 - d_{\ell,j})
\end{equation}
Under a Gaussian copula with correlation matrix $\mathbf{R}$:
\begin{equation}
D_j = 1 - C_{\mathbf{R}}(1 - d_{1,j}, \ldots, 1 - d_{L,j})
\end{equation}
At correlation $\rho = 1$, the second layer adds nothing ($D = d$). At $\rho = 0$, we recover independence.
\end{theorem}

Three layers at individually modest rates ($d_1 = 0.3$, $d_2 = 0.75$, $d_3 = 0.5$) yield $D = 1 - (0.7)(0.25)(0.5) = 0.9125$, exceeding 91\% cumulative DRE. Achieving DRE $> 95\%$ requires $\geq 3$ diverse techniques \citep{capersjones2012}. A weak independent detector ($d = 0.3, \rho = 0$) can outperform a strong correlated detector ($d = 0.7, \rho = 0.8$) in marginal DRE improvement.

\subsection{The Accretion Category}

\begin{definition}[Accretion Defect]
\label{def:accretion}
An \emph{accretion defect} is a code change that is: (i) functionally correct (all tests pass), (ii) superfluous (there exists a smaller change $\Delta' \subset \Delta$ satisfying the same requirement), (iii) individually defensible (it passes all automated checks and violates no stated standard), and (iv) collectively harmful (cumulative application increases codebase entropy by $\epsilon > 0$ per change).
\end{definition}

The Accretion Category is absent from all classical defect taxonomies (IEEE 1044, ODC, Beizer, CWE) because all four assume intentional authorship: code results from a human decision that was either correct or incorrect. Accretion arises from statistical pattern completion, which produces code because it is \emph{probable}, not because it is \emph{necessary}. It is the characteristic defect class of AI-generated code, and the primary driver of the GitClear observations (8$\times$ duplication, 60\% refactoring collapse).

Individual accretion detection is undecidable (requires exhibiting a smaller sufficient change, reducing to program equivalence). Aggregate detection is feasible via statistical process control on entropy proxies: accretion rate, refactoring ratio, clone frequency, lines per function point.

\begin{corollary}[Refactoring Ratio Threshold]
\label{cor:refactoring}
Let $r$ be the fraction of refactoring changes ($\alpha < 0$) and $(1-r)$ additions ($\alpha > 0$). Equilibrium requires $r \cdot |\mathbb{E}[\alpha | \text{refactoring}]| \geq (1-r) \cdot \mathbb{E}[\alpha | \text{addition}]$. GitClear data indicates the pre-AI equilibrium at $r \approx 0.24$; the post-AI state at $r \approx 0.095$ violates the bound by approximately 2.5$\times$.
\end{corollary}

\subsection{Cost-of-Quality Model}

The Quality pillar develops a cost-of-quality model following the Juran tradition:

\begin{equation}
    \text{CoQ} = C_{\text{prevention}} + C_{\text{appraisal}} + C_{\text{internal failure}} + C_{\text{external failure}}
\end{equation}

The classical 1:10:100 ratio (prevention to internal failure to external failure) requires recalibration by decidability class:

\begin{itemize}
\item \textbf{Decidable types:} approximately 1:3:100. Detection is cheap (automated at the appraisal stage); production costs remain high.
\item \textbf{Semi-decidable types:} approximately 1:15:100. Appraisal requires LLM or human review.
\item \textbf{Undecidable types:} approximately 1:50:30. Appraisal is expensive (deep expert review); but production cost is \emph{lower} because these degrade maintainability, not availability.
\end{itemize}

For Undecidable types, appraisal cost can exceed discounted external failure cost. It is sometimes economically rational to \emph{accept} these defects.

\begin{theorem}[Optimal Quality Level]
\label{thm:optquality}
With defect rate $\lambda(q) = \lambda_0 e^{-\beta q}$, the optimal quality investment is:
\begin{equation}
q^* = \frac{1}{\beta} \ln\!\left(\frac{c_F \cdot \lambda_0 \cdot \beta}{c_P + c_A}\right)
\end{equation}
where $c_F$ is failure cost and $c_P + c_A$ is marginal prevention/appraisal cost.
\end{theorem}

\begin{proof}[Proof sketch]
Total cost $\text{CoQ}(q) = (c_P + c_A) \cdot q + c_F \cdot \lambda_0 e^{-\beta q}$. Differentiating and setting to zero: $(c_P + c_A) = c_F \cdot \lambda_0 \cdot \beta \cdot e^{-\beta q^*}$. Solving for $q^*$ yields the stated expression. The second derivative is positive, confirming a minimum.
\end{proof}

This gives context-dependent quality targeting: $q^* \approx 1$ for security-critical production code ($c_F$ high), $q^* \approx 0.7$ for internal tools ($c_F$ moderate, Layers 1--2 suffice), $q^* \approx 0.3$ for prototypes ($c_F$ low, Layer 1 only). The appraisal compression problem further constrains the achievable quality: if code generation rate $g$ exceeds appraisal capacity $a$, a fraction $(g-a)/g$ bypasses appraisal entirely. Sonar data (48\% always verify) implies $g/a \approx 2$ on average; roughly half of AI code receives no appraisal \citep{sonar2026}. The quality stack's ordering is not merely cost optimization; it is a \emph{throughput constraint}: only automated layers (1--2) can match AI generation speed.


% ====================================================================
% 8. THE GOVERNANCE PILLAR (EXPANDED)
% ====================================================================
\section{The Governance Pillar}
\label{sec:governance}

The Governance pillar addresses the system-level problem: how do you maintain architectural coherence when AI agents modify codebases faster than humans can evaluate the implications?

\subsection{Theory Loss as Rate-Distortion}

Following \citet{naur1985}, a \emph{program theory} $T$ is the totality of a developer's understanding: the mapping from problem domain to code, design rationale, rejected alternatives, and the causal model of change propagation. The Governance pillar formalizes theory externalization as a rate-distortion problem:

\begin{theorem}[Theory Externalization Bound]
\label{thm:theory-loss}
The minimum description length of a program theory $T$ under distortion constraint $D$ satisfies:
\begin{equation}
    R(D) = \min_{p(\hat{T}|T): \mathbb{E}[d(T,\hat{T})] \leq D} I(T; \hat{T})
\end{equation}
where $\hat{T}$ is the externalized representation (documentation, ADRs, comments) and $d(T, \hat{T})$ is a distortion measure.
\end{theorem}

\begin{conjecture}[Tacit Divergence]
\label{conj:tacit}
Decompose the program theory as $T = (T_{\mathrm{explicit}}, T_{\mathrm{tacit}})$. Model $T_{\mathrm{explicit}}$ as a discrete random variable over a finite alphabet (design decisions, API choices, data structure selections) and $T_{\mathrm{tacit}}$ as a continuous random variable (intuitive weightings, aesthetic judgments, causal models of change propagation). Under squared-error distortion for the continuous component:
\begin{equation}
    \lim_{D \to 0} R_{\mathrm{tacit}}(D) = \infty
\end{equation}
with $R_{\mathrm{tacit}}(D) = \Omega(\log(1/D))$ as $D \to 0$. In contrast, $R_{\mathrm{explicit}}(D)$ remains finite for all $D \geq 0$ (since a discrete source with finite entropy has bounded rate-distortion).
\end{conjecture}

\begin{remark}
The conjecture's formal content depends on the modeling choice: treating tacit knowledge as a continuous source makes the divergence a standard result from rate-distortion theory (for a Gaussian source, $R(D) = \frac{1}{2}\log(\sigma^2/D)$, which diverges logarithmically). The substantive claim is that the continuous-source model is appropriate for tacit knowledge, which is a modeling assumption motivated by Polanyi's philosophical arguments rather than a mathematical fact. If tacit knowledge were instead a discrete source with finite entropy, the conjecture would be false. The conjecture is therefore testable in principle: demonstrate that design intuitions can be fully captured by a finite codebook, and the conjecture is falsified.
\end{remark}

If the Tacit Divergence Conjecture holds, then no governance scheme can fully externalize program theory. Some knowledge will always be lost when the original developers (human or AI) are no longer involved. This is a fundamental limitation, not an engineering challenge to be solved.

\subsection{Collins's Tacit Knowledge Taxonomy}

\citet{collins2010} advances beyond \citet{polanyi1966} by taxonomizing tacit knowledge into three types with distinct externalization properties:

\begin{center}
\begin{tabular}{lp{5cm}p{5cm}}
\toprule
Type & What Governance Can Capture & What Resists Capture \\
\midrule
RTK (Relational) & Decision context, rationale, rejected alternatives, advice received & Obvious-to-the-writer context \\
STK (Somatic) & Checklists, review heuristics, examples (lossy) & Debugging intuition, code quality ``smell'' \\
CTK (Collective) & (Nothing) & Team norms, ``good enough'' thresholds \\
\bottomrule
\end{tabular}
\end{center}

Relational tacit knowledge (RTK) is knowledge that \emph{could} be made explicit but has not been, due to social contingency; ADRs and governance tools operate effectively here. Somatic tacit knowledge (STK) resides in embodied practice, partially capturable through instrumentation but not transmissible through text. Collective tacit knowledge (CTK) is ``embodied in society,'' acquired only through enculturation, and resistant to externalization by any known means.

The Dreyfus model adds a further constraint: expert knowledge is specifically \emph{less} amenable to externalization than novice knowledge. The more expert a practitioner, the more their knowledge resides in pattern recognition that resists decomposition into rules. AI agents operating at the ``competent'' level of the Dreyfus hierarchy lack the intuitive, holistic perception of proficient and expert practitioners.

\subsection{The False Confidence Problem}

If \citet{naur1985} is correct, then any governance artifact (THEORY.md, structured ADRs, extracted decision logs) is a lossy compression. The map-territory problem applies: documentation is a map, not the territory. The risk: a team possessing documentation may believe it understands the system's theory when it actually possesses only a compressed, lossy representation. This could be worse than having no documentation, because no documentation produces appropriate humility (``we don't know why this is here''), while lossy documentation produces false confidence (``THEORY.md says it's for X, so we'll change it accordingly'').

However, this argument has practical limits. Lossy compression is still enormously useful (JPEG images demonstrate this). Documentation with explicit uncertainty markers (``confidence: medium,'' ``last validated: 2025-03'') addresses the false confidence risk.

\subsection{The Double Bottleneck Theorem}

AI coding agents produce verbose execution logs $L$ (often 100K+ tokens). A governance system must extract structured knowledge $(D, A, R)$ (decisions, assumptions, rejected alternatives). This forms a Markov chain: $T \to L \to (D, A, R) \to G$, where $G$ represents governance actions.

\begin{theorem}[Double Bottleneck]
\label{thm:double-bottleneck}
By the data processing inequality, applied twice:
\begin{equation}
    I(T; G) \leq I(T; (D, A, R)) \leq I(T; L)
\end{equation}
Each stage of the governance pipeline loses information. Governance fidelity is bounded above by log richness.
\end{theorem}

\begin{corollary}[Log Richness Lower Bound]
For governance to achieve fidelity $I(T; G) \geq F_{\min}$, the agent log must satisfy $I(T; L) \geq F_{\min}$. Terse logs impose an upper bound on governance fidelity regardless of extraction sophistication.
\end{corollary}

\subsection{The Governance Capacity Bound}

\begin{proposition}[Governance Capacity Principle]
\label{thm:governance-capacity}
Let $R_{\text{drift}}$ be the rate at which architectural drift is injected (bits per time unit) and $C_{\text{gov}}$ the governance channel capacity (bits per time unit the governance system can process and correct). Then:
\begin{enumerate}
    \item If $R_{\text{drift}} < C_{\text{gov}}$, there exists a governance coding scheme (combination of reviews, tests, policies, automated checks) maintaining architectural coherence with arbitrarily small residual drift.
    \item If $R_{\text{drift}} > C_{\text{gov}}$, \textbf{no governance scheme} can prevent architectural drift from growing without bound.
\end{enumerate}
\end{proposition}

\begin{proof}[Proof sketch]
By structural analogy to Shannon's noisy channel coding theorem. Part (1): when drift rate is below capacity, governance corrections can be encoded with sufficient redundancy. The ``coding scheme'' corresponds to layered governance (linting catches syntactic drift, CI tests catch behavioral drift, human review catches semantic drift, ADRs catch strategic drift). Part (2): by the converse of the channel coding theorem, when drift rate exceeds capacity, governance failure probability is bounded away from zero.
\end{proof}

\begin{remark}
This result is framed as a \emph{structural analogy} to Shannon's noisy channel coding theorem, not a direct application. Shannon's theorem requires a precisely defined channel with stationary transition probabilities and a finite input alphabet; the ``governance channel'' lacks these properties in the strict sense. The principle's value is in design guidance: it formalizes the intuition that governance systems have finite corrective bandwidth, that this bandwidth is measurable in principle, and that exceeding it produces qualitatively different (unbounded) behavior.
\end{remark}

For a layered governance system with $L$ independent layers, the combined capacity satisfies $C_{\text{gov}} \leq \sum_{\ell=1}^{L} C_\ell$, with equality when layers operate on non-overlapping aspects. When AI agents increase the modification rate $\lambda$ while $C_{\text{gov}}$ remains constant (human review bandwidth is fixed), the condition $R_{\text{drift}} > C_{\text{gov}}$ is eventually satisfied. Automated governance layers (fitness functions, architectural tests, invariant checkers) are therefore \emph{necessary}, not merely convenient.

\subsection{Multi-Granularity Governance as Cascade Control}

The Governance pillar models the governance system as a cascade control system with three loops:

\begin{enumerate}
    \item \textbf{Inner loop} (synchronous, per-generation): validates each code generation against syntax rules, linting constraints, and file-ownership invariants. Sampling period $T_1$ (fastest). Bandwidth: high.
    \item \textbf{Middle loop} (asynchronous, per-phase): runs compliance audits, reflexion model checks, and contract verification at phase boundaries. Sampling period $T_2 \gg T_1$. Bandwidth: moderate.
    \item \textbf{Outer loop} (deliberative, per-milestone): performs ATAM tradeoff analysis, strategic alignment review, and ADR lifecycle management. Sampling period $T_3 \gg T_2$. Bandwidth: low.
\end{enumerate}

Let $D(z)$ denote drift introduced by agent-generated code. Each loop has controller $C_k(z)$ and plant $P_k(z)$. The inner loop's closed-loop transfer function is:
\begin{equation}
    G_{\mathrm{inner}}(z) = \frac{C_1(z) P_1(z)}{1 + C_1(z) P_1(z)}
\end{equation}
Each outer loop treats the inner loop's output as its plant.

\begin{theorem}[Cascade Governance Stability]
\label{thm:cascade-stability}
The three-loop system is stable if and only if:
\begin{enumerate}
    \item Each loop is individually stable (all poles of $G_k(z)$ lie within the unit circle).
    \item Bandwidth separation holds: $\omega_{\mathrm{bw},1} > \omega_{\mathrm{bw},2} > \omega_{\mathrm{bw},3}$.
    \item Each outer loop treats the inner as approximately unity gain within its bandwidth.
\end{enumerate}
\end{theorem}

\begin{proof}[Proof sketch]
Condition (1) follows from BIBO stability of discrete-time systems. Condition (2) ensures frequency-domain decoupling; without it, outer loop corrections feed into the inner loop's operating range, creating destructive interference. Condition (3) ensures the outer loop's model of the inner loop is accurate within its bandwidth. Under these conditions, the combined characteristic polynomial factors approximately into three independent polynomials.
\end{proof}

Classical cascade control requires a 3:1 to 5:1 bandwidth ratio between adjacent loops. If the inner loop checks every generation (seconds), the middle loop should operate on minutes-to-hours timescales, and the outer loop on days-to-weeks.

\subsection{Nyquist Sampling Requirements}

Let $\omega_d^{(k)}$ be the maximum frequency of architectural drift relevant to loop $k$. Each governance loop must sample at:
\begin{equation}
    f_k \geq 2\omega_d^{(k)}
\end{equation}

Concrete rate calculations illustrate the constraint. If an AI agent produces 50 code generations per hour, the inner loop must sample at least 100 events per hour to capture per-generation violations. If emergent drift (coupling growth, convention erosion) manifests on a per-PR timescale and a team merges 10 PRs per day, the middle loop requires at least 20 audit events per day. If strategic misalignment accumulates over weeks, the outer loop requires review at least every 3.5 days to satisfy the Nyquist condition.

If the outer loop samples too slowly, it aliases gradual architectural erosion as apparent stability, missing drift until it becomes catastrophic. When AI agents increase the disturbance frequency (more changes per unit time), the sampling rate of each loop must increase accordingly, or the loop becomes unstable.

\subsection{The Governance Bifurcation}

Define architectural coherence $C(t) \in [0,1]$ as a scalar conformance measure. The dynamics are:
\begin{equation}
    \frac{dC}{dt} = \mu \cdot G(t) \cdot (1 - C)^\beta - \gamma \cdot R(t) \cdot (1 - C)
\end{equation}
where $G(t)$ is governance intervention rate, $R(t)$ is agent generation rate, $\mu$ is governance effectiveness, $\gamma$ is agent drift propensity (using $\gamma$ for drift propensity, not $\delta_a$, to avoid overloading the relaxation operator symbol), and $\beta > 1$ captures diminishing governance returns. The asymmetry ($\beta$ in the governance term but not the drift term) reflects the empirical observation that easy violations are caught first while subtle ones require disproportionate effort, whereas drift injection is approximately linear in agent activity.

\begin{theorem}[Governance Bifurcation]
\label{thm:bifurcation}
The system undergoes a transcritical bifurcation at $\mu G / \gamma R = 1$:
\begin{itemize}
    \item \textbf{Supercritical} ($\mu G > \gamma R$): a stable interior fixed point at $C^* = 1 - (\gamma R / \mu G)^{1/(\beta-1)} > 0$. The system self-corrects.
    \item \textbf{Subcritical} ($\mu G < \gamma R$): the only stable fixed point is $C^* = 0$ (total coherence collapse).
\end{itemize}
Near the bifurcation, recovery time diverges (critical slowing down), and small perturbations in $R(t)$ cause qualitative regime changes.
\end{theorem}

\begin{proof}[Proof sketch]
Setting $dC/dt = 0$ and factoring yields $(1 - C)[\mu G (1 - C)^{\beta - 1} - \gamma R] = 0$. The boundary fixed point $C = 1$ is always present and is stable (linearization gives $dC/dt \approx -\gamma R \cdot \Delta C$ near $C = 1$ since $(1-C)^{\beta-1} \to 0$ for $\beta > 1$), but it is unreachable from any $C < 1$ in finite time when $\beta > 1$, because both drift and governance forces vanish as $C \to 1$. The interior fixed point $C^* = 1 - (\gamma R / \mu G)^{1/(\beta - 1)}$ exists only when $\mu G > \gamma R$ and is the practically relevant attractor. When $\mu G < \gamma R$, the drift term dominates everywhere in $(0, 1)$, driving coherence to zero. The interior fixed point exchanges stability with $C = 0$ at $\mu G / \gamma R = 1$, producing a transcritical bifurcation at that critical ratio.
\end{proof}

The phase portrait has two regimes separated by a sharp transition. In the supercritical regime, perturbations away from $C^*$ are corrected by governance feedback. In the subcritical regime, coherence degrades monotonically. With $\beta > 1$, the equilibrium coherence is always strictly less than 1: perfect conformance is asymptotically unattainable, though the gap shrinks as $G/R$ increases.

\subsection{The Ratchet Lattice}

Let $\mathcal{L} = (2^{\mathcal{R}}, \subseteq)$ be the power-set lattice over a universe of governance rules $\mathcal{R}$. A governance state at time $t$ is an element $S_t \in \mathcal{L}$.

\begin{definition}[Governance Ratchet]
\label{def:ratchet}
A ratchet is a closure operator $\rho: \mathcal{L} \to \mathcal{L}$ on the constraint lattice $\mathcal{L}$, satisfying: (i) extensiveness: $S \subseteq \rho(S)$ (rules are only added); (ii) monotonicity: $S \subseteq S' \implies \rho(S) \subseteq \rho(S')$; (iii) idempotence: $\rho(\rho(S)) = \rho(S)$.
\end{definition}

The ADR acceptance workflow is literally a closure operation: accepting a decision adds constraints, those constraints may force additional constraints (transitive closure), and the result is stable under re-application.

\begin{theorem}[Ratchet Convergence]
\label{thm:ratchet-convergence}
Let $\rho$ be a governance ratchet on a finite rule universe $\mathcal{R}$. The sequence $S_0, S_1 = \rho(S_0), S_2 = \rho(S_1), \ldots$ converges to a fixed point $S^* = \rho(S^*)$ in at most $|\mathcal{R}|$ steps.
\end{theorem}

\begin{proof}[Proof sketch]
By extensiveness, $S_t \subseteq S_{t+1}$, so the sequence is monotonically non-decreasing. Since $\mathcal{R}$ is finite, $\mathcal{L}$ has finite height $|\mathcal{R}|$. A monotone chain of finite height stabilizes. Alternatively, by the Knaster-Tarski theorem, every monotone operator on a complete lattice has a least fixed point.
\end{proof}

\begin{theorem}[Ossification Risk]
\label{thm:ossification}
If $\mathcal{R}$ contains contradictory rules ($\exists\, r_1, r_2 \in \mathcal{R}$ with no code satisfying both) and the ratchet is purely additive, then any trajectory including both $r_1$ and $r_2$ produces an ossified state (where $\mathrm{Allowed}(S) = \emptyset$).
\end{theorem}

\begin{proof}
By extensiveness, once $r_1 \in S_t$ and $r_2 \in S_{t'}$, both persist in all subsequent states. If $r_1 \wedge r_2$ is unsatisfiable, $\mathrm{Allowed}(S) = \emptyset$.
\end{proof}

\subsection{Controlled Relaxation with Evidence Expiry}

The ossification theorem provides formal justification for relaxation mechanisms. Define a relaxation operator $\delta: \mathcal{L} \to \mathcal{L}$ with $\delta(S) \subseteq S$. The combined governance operator is:
\begin{equation}
    \Gamma(S) = \rho(S \setminus \delta(S))
\end{equation}
First remove expired or superseded rules via $\delta$, then apply the ratchet closure $\rho$. Relaxation may be triggered by evidence expiry (a rule's justifying evidence has a validity window), decision supersession (a newer ADR replaces an older one), or explicit deprecation.

\begin{theorem}[Controlled Descent]
If $\delta$ removes only rules whose justifying evidence has expired, and $\rho$ re-derives still-justified consequences, then $\Gamma$ produces states that are both internally consistent and evidence-backed.
\end{theorem}

The evidence expiry window for each rule should be proportional to its domain-specific median lifetime, connecting the ratchet to the survival analysis below.

\subsection{Decision Survival Analysis with Competing Risks}

Architectural decisions have finite lifetimes. The survival function is $S(t) = P(\text{decision valid at time } t) = \exp\left(-\int_0^t h(u) \, du\right)$ where $h(u)$ is the hazard function. Decisions fail for multiple independent reasons. Let $K$ index competing risks:

\begin{center}
\begin{tabular}{clp{7cm}}
\toprule
$k$ & Risk Type & Description \\
\midrule
1 & Technology obsolescence & Chosen technology superseded \\
2 & Requirement change & Business requirements shift \\
3 & Evidence expiry & Justifying data becomes stale \\
4 & Organizational change & Team structure or strategy shifts \\
\bottomrule
\end{tabular}
\end{center}

The cause-specific hazard for risk $k$ is:
\begin{equation}
    h_k(t) = \lim_{\Delta t \to 0} \frac{P(t \leq T_d < t + \Delta t, \; K_d = k \mid T_d \geq t)}{\Delta t}
\end{equation}
with overall hazard $h(t) = \sum_k h_k(t)$. The cumulative incidence function is:
\begin{equation}
    F_k(t) = \int_0^t h_k(u) \exp\!\left(-\int_0^u h(\tau)\, d\tau\right) du
\end{equation}

Empirically calibrated decay rates (using Weibull hazard $h(t) = (k/\lambda)(t/\lambda)^{k-1}$ with shape $k$ and scale $\lambda$):

\begin{center}
\begin{tabular}{lcccc}
\toprule
Decision Domain & Dominant Risk & Weibull $\alpha$ & Scale $\lambda$ (mo.) & Median Life \\
\midrule
Frontend framework & $k=1$ & 1.5--2.5 & 2--5 & 1--4 mo. \\
API contract & $k=2$ & 1.2--1.8 & 5--12 & 3--9 mo. \\
Database schema & $k=3$ & 1.0--1.4 & 18--36 & 12--36 mo. \\
Security policy & $k=4$ & 0.8--1.2 & 12--24 & 8--18 mo. \\
\bottomrule
\end{tabular}
\end{center}

These parameter ranges are \emph{illustrative estimates} derived from practitioner experience and ecosystem churn rates. No controlled longitudinal study of decision decay rates exists. The qualitative ordering (frontend decisions decay fastest; database decisions are most durable) is better supported than the specific numerical values.

\subsection{The Theory Preservation Problem}

The Governance pillar identifies the \emph{theory preservation problem} as fundamentally open: how do you preserve program theory across agent generations, context window boundaries, and team transitions? This is not merely a documentation problem. Documentation captures the explicit component of theory; the tacit component (Naur's ``direct, intuitive knowledge,'' Polanyi's personal knowledge, Collins's mimeomorphic actions) resists formalization.

When AI generates a substantial fraction of code, a new failure mode appears: \emph{theoretically orphaned implementations}, code for which the theory was never in anyone's head. Current approaches (Architecture Decision Records, ARCHITECTURE.md, embedded comments, convention-by-example) are partial solutions. Each captures some explicit knowledge; none captures tacit knowledge. The theory preservation problem is the long-term challenge that will determine whether AI-assisted software development produces maintainable systems or disposable ones.
% ====================================================================
% EXPANDED SECTIONS 9, 10, 12, AND APPENDIX A
% For: Towards a Science for AI Coding Agent Harnesses
% ====================================================================
% This file contains expanded, self-contained versions of:
%   Section 9:  Cross-Pillar Synthesis
%   Section 10: Empirical Foundations
%   Section 12: Limitations and Open Problems
%   Appendix A: Cross-Pillar Theorem Index
% ====================================================================


% ====================================================================
% 9. CROSS-PILLAR SYNTHESIS
% ====================================================================
\section{Cross-Pillar Synthesis}
\label{sec:synthesis}

The preceding sections presented each pillar independently. This section, the primary contribution of this synthesis paper, identifies the structural connections that bind them into a unified framework.

\subsection{Information as the Unifying Currency}

The deepest structural connection across all seven pillars is that each reasons about \emph{information flow} through the harness:

\begin{itemize}[nosep]
    \item \textbf{Abstraction}: the information gap $K(P \mid S)$ between specification and code.
    \item \textbf{Information}: the Shannon decomposition $H(P \mid S) = H(P \mid S,C,\theta) + I(P;C \mid S,\theta) + I(P;\theta \mid S)$.
    \item \textbf{Reliability}: compound error as information destruction, $R(n) = p^n$.
    \item \textbf{Coordination}: error amplification as correlated information loss across agents.
    \item \textbf{Temporal}: context fidelity $\Phi(t) = e^{-\Lambda t}$ as information decay over time.
    \item \textbf{Quality}: the Detection Ceiling $d_{\ell,j} \leq I(X; D_j)/H(D_j)$ as an information bound on verification.
    \item \textbf{Governance}: governance capacity $C_{\text{gov}}$ as information throughput of the correction channel.
\end{itemize}

The unification is strongest along the \emph{semantic axis} (Abstraction, Information) and the \emph{assurance axis} (Quality, Governance), where information theory is the native formalism. For the \emph{execution axis} (Reliability, Coordination, Temporal), the information-theoretic framing is a valid interpretation but not the generative formalism: $R(n) = p^n$ is a probability product from reliability theory, the Extended Amdahl's Law draws on parallel computing and concurrency control, and VIH draws on scheduling and queueing theory. Calling compound errors ``information destruction'' is accurate (errors do corrupt the information content of pipeline output) but retrospective; the results would stand unchanged without the information-theoretic gloss. We adopt the information lens because it provides a common vocabulary for cross-pillar reasoning, while acknowledging that its explanatory power varies across pillars.

\subsection{The Abstraction Gap Chain}

The first structural connection links three pillars in a chain:

\begin{enumerate}
    \item \textbf{Abstraction} defines the gap: $\mathcal{G}(S,P) = K(P \mid S)$.
    \item \textbf{Information} operationalizes it: $H(P \mid S) = H(P \mid S,C,\theta) + I(P;C \mid S,\theta) + I(P;\theta \mid S)$, decomposing the gap into three actionable levers.
    \item \textbf{Quality} addresses what happens when the gap is traversed poorly: the 18 slop types are failure modes of the translation from spec to code, and the Compound Degradation Theorem (Theorem~\ref{thm:compound-degradation}) quantifies the consequences.
\end{enumerate}

The chain makes precise how a specification becomes code and what goes wrong along the way. The Abstraction pillar establishes the theoretical limits. The Information pillar identifies the levers. The Quality pillar classifies the failure modes and designs defenses.

\subsection{The Compound Error Cascade}

The second structural connection links three pillars through the theme of multiplicative degradation:

\begin{enumerate}
    \item \textbf{Reliability} proves $R(n) = p^n$: compound errors in sequential pipelines.
    \item \textbf{Coordination} extends to multi-agent settings: the Extended Amdahl's Law (Theorem~\ref{thm:extended-amdahl}) adds the reliability factor $(1-\delta_a)^n$, and empirical amplification ($17.2\times$) exceeds the independence prediction.
    \item \textbf{Temporal} frames the rate at which compound errors occur: VIH $= \lambda_{\text{raw}} \cdot p_{\text{verify}}$ captures the tradeoff between iteration speed and error rate.
\end{enumerate}

The cascade reveals a fundamental tension: the same exponential that makes compound errors devastating also makes per-step improvements highly leveraged (elasticity $= n$). This is the core design insight: invest in per-step quality, not pipeline sophistication.

\subsection{The Structural Enforcement Principle}

The principle that structural enforcement dominates instructional enforcement appears independently in four pillars:

\begin{table}[h]
\centering
\caption{The structural enforcement principle across pillars.}
\label{tab:structural}
\begin{tabular}{lp{10cm}}
\toprule
\textbf{Pillar} & \textbf{Manifestation} \\
\midrule
Information & Tier boundaries enforced structurally (file system, .gitignore) vs. instructionally (prompt rules). Compliance gap: $(1-\epsilon)^T$. \\
Reliability & Cost-of-failure threshold $L^*$ determines when structural enforcement is cost-effective. Threshold drops as $1/n$ for multi-step pipelines. \\
Quality & Decidable slop types should be caught by structural gates (Layer 1), not by probabilistic methods. \\
Coordination & File-ownership contracts enforce $F_i \cap F_j = \emptyset$ structurally, eliminating textual/syntactic conflicts by construction. \\
\bottomrule
\end{tabular}
\end{table}

The convergence is not coincidental. Structural enforcement converts probabilistic guarantees into deterministic ones. In any system where errors compound (as they do in all agent pipelines), converting even a small fraction of probabilistic steps to deterministic ones has outsized impact on end-to-end reliability.

\subsection{The Four-Layer Defense Architecture}

The four-layer verification architecture appears in both the Reliability and Quality pillars:

\begin{table}[h]
\centering
\caption{The shared four-layer verification architecture.}
\label{tab:four-layer}
\begin{tabular}{clccc}
\toprule
\textbf{Layer} & \textbf{Mechanism} & \textbf{Cost} & \textbf{Determinism} & \textbf{Decidability Class} \\
\midrule
1 & Structural gates & Very low & Deterministic & Decidable \\
2 & Automated testing & Low & Probabilistic & Semi-decidable \\
3 & AI review & Moderate & Probabilistic & Semi-decidable \\
4 & Human review & High & Judgment & Undecidable \\
\bottomrule
\end{tabular}
\end{table}

The ordering by cost and determinism is not arbitrary; it follows from the decidability classification. Decidable defects should be caught at Layer 1 (structural gates), where detection is certain and marginal cost is near zero. Semi-decidable defects require Layers 2--3, where detection is probabilistic but cost-effective. Undecidable defects require Layer 4, where human judgment provides the irreplaceable element.

\subsection{Governance Capacity as the Macro Constraint}

The Governance pillar's capacity bound (Theorem~\ref{thm:governance-capacity}) is the system-level analog of the Reliability pillar's compound error sensitivity (Theorem~\ref{thm:compound-sensitivity}):

\begin{itemize}[nosep]
    \item \textbf{Reliability}: at the step level, if $p < 1$, errors compound as $p^n$.
    \item \textbf{Governance}: at the system level, if $R_{\text{drift}} > C_{\text{gov}}$, divergence accumulates without bound.
\end{itemize}

Both are threshold phenomena: below the threshold, the system is self-correcting; above it, degradation is unbounded. Both have the same structural form: a rate of degradation versus a capacity for correction. And both yield the same design prescription: invest in increasing the correction capacity (per-step quality, governance bandwidth) rather than adding post-hoc repair mechanisms.

\subsection{Shared Formal Machinery}
\label{sec:shared-machinery}

Table~\ref{tab:machinery} maps the 16 distinct formal techniques used across the seven pillars. The density of the mapping (most techniques appear in 2+ pillars) reflects the deep structural unity of the framework; the pillars are not independent theories glued together but facets of a common problem viewed through different mathematical lenses.

\begin{table}[h]
\centering
\caption{Formal machinery shared across pillars. A bullet (\textbullet) indicates that the pillar makes substantive use of the technique; a circle ($\circ$) indicates secondary or analogical use.}
\label{tab:machinery}
\small
\begin{tabular}{l@{\hskip 4pt}c@{\hskip 4pt}c@{\hskip 4pt}c@{\hskip 4pt}c@{\hskip 4pt}c@{\hskip 4pt}c@{\hskip 4pt}c}
\toprule
\textbf{Technique} & \textbf{Abs} & \textbf{Inf} & \textbf{Rel} & \textbf{Crd} & \textbf{Tmp} & \textbf{Qlt} & \textbf{Gov} \\
\midrule
Shannon information theory ($H$, $I$, DPI) & \textbullet & \textbullet & $\circ$ & $\circ$ & $\circ$ & \textbullet & \textbullet \\
Kolmogorov complexity ($K$) & \textbullet & $\circ$ & & & & & \textbullet \\
Series reliability / compound errors ($p^n$) & & & \textbullet & \textbullet & $\circ$ & & \\
Structural vs. instructional enforcement & & \textbullet & \textbullet & \textbullet & & \textbullet & $\circ$ \\
Submodular optimization / greedy approx. & & \textbullet & & & & \textbullet & \\
Control theory (cascade, Nyquist, stability) & & & & & $\circ$ & & \textbullet \\
Lattice theory (Galois, closure, refinement) & \textbullet & & & & & & \textbullet \\
Concurrency control / isolation levels & & & & \textbullet & & & \\
Survival analysis (competing risks, Weibull) & & & & & & & \textbullet \\
Scheduling theory (Amdahl, stochastic DAG) & & & & \textbullet & \textbullet & & \\
FKG / correlation inequalities & & & \textbullet & & & \textbullet & \\
Graph partitioning (min-cut, Cheeger) & & & & \textbullet & & & \\
Rate-distortion theory & & & & & & & \textbullet \\
Queueing theory (Little's law, Kingman) & & & & & \textbullet & & \\
Speculative execution theory & & & & & \textbullet & & \\
Cache competitive analysis (Sleator-Tarjan) & & & & & \textbullet & & \\
\bottomrule
\end{tabular}
\end{table}

Three patterns emerge. First, Shannon information theory provides the broadest cross-pillar vocabulary, appearing substantively in four pillars and secondarily in three more. Second, structural enforcement is the most widely shared \emph{design principle} (four pillars derive it independently). Third, several techniques (concurrency control, survival analysis, cache competitive analysis) are used in only one pillar, suggesting that further cross-pollination may be productive; for instance, survival analysis could model the lifetime of context relevance (connecting Governance to Temporal), and concurrency control formalisms could sharpen the governance ratchet's interaction with parallel agent modifications.

\subsection{Shared Empirical Sources}
\label{sec:shared-empirical}

Table~\ref{tab:empirical-sources} maps the 14 primary empirical sources to the pillars that rely on them. The density of cross-pillar citations highlights both a strength (convergent evidence from multiple analytical perspectives) and a vulnerability (the same data points bear disproportionate evidentiary weight).

\begin{table}[h]
\centering
\caption{Key empirical sources and the pillars that draw on them. Columns indicate pillar usage: Abs(traction), Inf(ormation), Rel(iability), Crd (Coordination), Tmp (Temporal), Qlt (Quality), Gov(ernance).}
\label{tab:empirical-sources}
\small
\begin{tabular}{l@{\hskip 3pt}c@{\hskip 3pt}c@{\hskip 3pt}c@{\hskip 3pt}c@{\hskip 3pt}c@{\hskip 3pt}c@{\hskip 3pt}c@{\hskip 6pt}p{4.5cm}}
\toprule
\textbf{Source} & \textbf{Abs} & \textbf{Inf} & \textbf{Rel} & \textbf{Crd} & \textbf{Tmp} & \textbf{Qlt} & \textbf{Gov} & \textbf{Key Finding} \\
\midrule
GitClear 2025 & & \textbullet & & & \textbullet & \textbullet & \textbullet & 211M lines; 8$\times$ duplication; 60\% refactoring collapse \\
DORA 2025 & & & & \textbullet & \textbullet & \textbullet & \textbullet & Throughput up 21\%; stability negative \\
Kim et al.\ 2025 & \textbullet & & & \textbullet & & & & $17.2\times$ error amplification; 180 configs \\
METR 2026 & \textbullet & & \textbullet & & & \textbullet & & 50\% test-passing PRs unmerged \\
Liu et al.\ 2024 & & \textbullet & & & & & & Lost in the Middle; U-shaped recall \\
Chroma 2025 & & \textbullet & & & \textbullet & & & 18 models; monotonic degradation \\
JetBrains 2025 & & \textbullet & & & & & & Complexity increase in AI code \\
CodeRabbit 2025 & & & & & & \textbullet & & 1.68$\times$ issues; 2.74$\times$ XSS vulns \\
Veracode 2025 & & & & & & \textbullet & & AI code security defect rates \\
Apollo Research 2025 & & & \textbullet & & & \textbullet & & Scheming behaviors in agents \\
NIST CAISI 2025 & & & \textbullet & & & \textbullet & & Benchmark cheating evidence \\
Hariharan et al.\ 2025 & & & \textbullet & & & & & 63\% failure on 100-step tasks \\
Spotify Honk 2025 & & \textbullet & & & & \textbullet & & Production harness telemetry \\
Gloaguen et al.\ 2026 & & \textbullet & & & & & & AGENTS.md convention data \\
\bottomrule
\end{tabular}
\end{table}

The table reveals a concentration risk: GitClear, DORA, and METR each support claims in 3--4 pillars, meaning that a successful challenge to any one of these sources would ripple across the framework. Section~\ref{sec:empirical} provides a quality assessment of each source to make this risk explicit.

\subsection{Additional Cross-Pillar Connections}
\label{sec:missed-connections}

Beyond the three primary cascades (abstraction gap chain, compound error cascade, structural enforcement principle), five cross-pillar interactions merit formal development.

\paragraph{1. Accretion meets the ratchet.}
The Quality pillar's Accretion Category (individually CI-passing changes that collectively degrade codebase health) is precisely the failure mode that the Governance pillar's ratchet mechanism (Definition~\ref{def:ratchet}) aims to prevent. The interaction, however, is non-trivial. Accretion defects are correct by construction; they satisfy all existing ratchet constraints. To prevent accretion, the ratchet must enforce \emph{structural health metrics} (duplication rate, dependency fan-out, abstraction depth) rather than per-change correctness. But structural health constraints are inherently more brittle than correctness constraints: they have higher false-positive rates and are sensitive to legitimate refactoring. Adding ratchet constraints on structural metrics risks triggering the ossification pathology (the Governance pillar's Ossification Theorem), where the constraint set grows until the allowed change set becomes empty or effectively so. The calibration question is: at what threshold should structural health metrics be ratcheted, and how should evidence expiry windows be set relative to the accretion rate $\alpha(t)$? The Compound Degradation Theorem (Theorem~\ref{thm:compound-degradation}) provides the growth rate; the Ratchet Convergence theorem provides the stabilization guarantee; what is missing is the bridging analysis showing how to set ratchet parameters that prevent accretion without inducing ossification.

\paragraph{2. Staleness compounds with coordination.}
The Temporal pillar's context fidelity decay $\Phi(t) = e^{-\Lambda(t-t_0)}$ interacts multiplicatively with the Coordination pillar's error amplification. In multi-agent settings, the effective decay rate $\Lambda$ is not a constant; it increases with the number of concurrent agents, because more agents means faster codebase evolution and therefore faster context invalidation. Specifically, if $k$ agents each commit at rate $\lambda_c$, the effective decay rate for any one agent's context is $\Lambda_{\text{eff}} \approx (k-1)\lambda_c \alpha$, where $\alpha$ is the fraction of commits that modify files in the agent's active context. This creates a feedback loop: more agents $\to$ faster staleness $\to$ higher per-agent error rate $\to$ more rework $\to$ more commits $\to$ even faster staleness. The Extended Amdahl's Law (Theorem~\ref{thm:extended-amdahl}) does not account for this feedback; incorporating staleness-induced error rate inflation into the reliability factor $(1-\delta_a)^n$ would likely reduce $n^*$ below the current estimate $1/\sqrt{\delta_a}$, particularly for tightly coupled codebases where $\alpha$ is large.

\paragraph{3. Lattice parallels across Abstraction and Governance.}
Both the Abstraction pillar and the Governance pillar employ lattice-theoretic structures, but with different orientations. The Abstraction pillar's refinement lattice orders specifications from most abstract (top, natural language) to most concrete (bottom, executable code), with the Galois connection $(\alpha, \gamma)$ mediating between specifications and programs. The Governance pillar's constraint lattice orders constraint sets by inclusion, with the ratchet as a closure operator that moves the system toward the join (more constrained). These two lattices are not independent: every specification in the refinement lattice \emph{induces} a constraint in the governance lattice (the specification constrains which programs are acceptable). A formal bridge would establish that the Galois connection on the refinement lattice is compatible with the closure operator on the constraint lattice; that is, refining a specification (descending the refinement lattice) should tighten constraints (ascending the constraint lattice) monotonically. If this compatibility holds, then the act of specification refinement automatically generates governance constraints, providing a formal foundation for ``specification as governance.''

\paragraph{4. Decomposition enables reuse discovery.}
The Coordination pillar's task decomposition (balanced min-cut on the coupling graph) and the Information pillar's reuse discovery problem are connected through the structure of the coupling graph itself. Balanced min-cut partitions identify clusters of highly coupled files; these clusters are also the natural units for reuse discovery, because semantic similarity concentrates within clusters. An agent assigned to a cluster via the decomposition has, by construction, the most relevant context for discovering reuse opportunities within that cluster. This suggests that coordination topology should be co-designed with the context architecture: the partition that minimizes merge conflicts (Coordination) simultaneously maximizes reuse-relevant context concentration (Information). The formal connection is that the coupling graph's cluster structure determines both the optimal decomposition and the natural reuse neighborhoods, and a single spectral analysis of the coupling graph can serve both objectives.

\paragraph{5. FKG amplification across coordination boundaries.}
The Quality pillar's FKG inequality (Proposition~\ref{prop:fkg}), which establishes that same-family LLM judges underestimate escape probabilities due to correlated blind spots, has a multi-agent analog that the current framework does not address. When $k$ agents from the same model family work on a shared codebase, their errors are positively correlated not only within each agent's pipeline (the Quality pillar's concern) but \emph{across} agents (the Coordination pillar's concern). The FKG inequality implies:
\begin{equation}
    P\!\left(\bigcap_{i=1}^k E_i\right) \geq \prod_{i=1}^k P(E_i)
\end{equation}
where $E_i$ is the event that agent $i$ introduces a defect in the shared semantic category. This means the Coordination pillar's independence assumption ($A_e \to k$) not only underestimates amplification (as the Kim et al.\ empirical data shows) but does so in a structurally predictable way: the FKG framework provides an \emph{upper} bound on the independence gap. Combining the FKG inequality with the coordination error amplification model would yield tighter bounds on multi-agent defect rates, explaining part of the $6\times$ discrepancy between the theoretical independence prediction and the empirical $17.2\times$ amplification.

\subsection{The Human Correction Cycle}
\label{sec:human-loop}

A significant limitation of the current framework is that it models agent pipelines as feedforward (Reliability) or parallel (Coordination) systems, with verification as a separate stage. In production harnesses, the dominant quality mechanism is the \emph{human correction cycle}: a user reads the agent's output, identifies where it went wrong, provides a correction or redirect, and the agent revises. This cycle operates continuously, not as a final gate.

The four-layer verification architecture (Section~\ref{sec:reliability}) places human review at Layer 4 (highest quality, highest cost, lowest throughput). This framing is technically correct in terms of cost per defect detected, but it misrepresents the role of humans in practice. Humans are not the last layer of a cascade; they are active co-pilots who steer the agent at every step. Human review is the only verification mechanism that can: (a) catch novel error categories not anticipated by automated layers, (b) redirect the agent when it is solving the wrong problem entirely, and (c) supply the tacit knowledge that the Governance pillar (Section~\ref{sec:governance}) identifies as unexternalizable.

The efficiency metric $\eta = d/c$ (detection rate per unit cost) penalizes human review for being expensive but does not credit it for handling the tail of the error distribution, including errors that no automated layer can detect. A complete formal treatment would model the human as a feedback channel with its own capacity constraints, latency characteristics, and learning dynamics (Bainbridge's Ironies apply: as automation improves, human monitors get less practice, making their interventions less reliable precisely when they are most needed). We leave this formalization as an important direction for future work and note that the design principles in Section~\ref{sec:principles} should be interpreted in the context of a system where human correction is the dominant quality lever, not merely the most expensive one.

\subsection{The Governance Information Budget}
\label{sec:gib}

We propose the \emph{Governance Information Budget} as the integrating concept across all seven pillars. The idea is that every harness operates under a total information budget, and the pillars decompose how that budget is allocated:

\begin{equation}
\label{eq:gib}
    \underbrace{K(P \mid S)}_{\text{Abstraction gap}} = \underbrace{I(P; C \mid S, \theta)}_{\text{Context (Information)}} + \underbrace{I(P; \theta \mid S)}_{\text{Prior (Model)}} + \underbrace{H(P \mid S, C, \theta)}_{\text{Residual (Harness must handle)}}
\end{equation}

The residual term $H(P \mid S, C, \theta)$ is the information that neither the context nor the model provides. This residual must be filled by the agent's exploration and reasoning, which is subject to:
\begin{itemize}[nosep]
    \item Compound error sensitivity (Reliability): each exploratory step has failure probability $q$, compounding as $(1-q)^n$.
    \item Error amplification (Coordination): multi-agent exploration amplifies errors by up to $17.2\times$.
    \item Context staleness (Temporal): the residual grows as context decays, $\Phi(t) = e^{-\Lambda t}$.
    \item Detection limits (Quality): some errors in filling the residual are undetectable, $d_{\ell,j} \leq I(X; D_j)/H(D_j)$.
    \item Governance capacity (Governance): the total rate at which the residual is filled must not exceed governance capacity $C_{\text{gov}}$.
\end{itemize}

\begin{proposition}[Governance Information Budget Constraint]
\label{prop:gib-constraint}
For a harness operating at steady state with change rate $\lambda$ (changes per unit time), the Governance Information Budget imposes five simultaneous constraints, one from each execution-axis and assurance-axis pillar:
\begin{align}
    \text{(Reliability)} &\quad \lambda \leq \frac{-\ln R_{\min}}{n \cdot |\ln p|} \label{eq:gib-rel} \\
    \text{(Coordination)} &\quad k \leq n^* \approx 1/\sqrt{\delta_a} \label{eq:gib-coord} \\
    \text{(Temporal)} &\quad \lambda \leq \Lambda / \ln(1/\Phi_{\min}) \label{eq:gib-temp} \\
    \text{(Quality)} &\quad \lambda \cdot P_{\text{esc}} \leq \mu_{\text{repair}} \label{eq:gib-qual} \\
    \text{(Governance)} &\quad \lambda \cdot H_{\text{drift/change}} \leq C_{\text{gov}} \label{eq:gib-gov}
\end{align}
where $R_{\min}$ is the minimum acceptable pipeline reliability, $\Phi_{\min}$ is the minimum acceptable context fidelity, $P_{\text{esc}}$ is the escape probability past all defense layers, $\mu_{\text{repair}}$ is the defect repair rate, and $H_{\text{drift/change}}$ is the architectural drift per change (in bits). The binding constraint (the one that limits $\lambda$ most severely) determines which pillar is the system's current bottleneck.
\end{proposition}

\begin{remark}
The five constraints in Proposition~\ref{prop:gib-constraint} are not independent. Increasing $k$ (more parallel agents) relaxes the Reliability constraint (more throughput) but tightens the Temporal constraint (faster staleness) and the Governance constraint (more drift per unit time). Increasing per-step quality $p$ relaxes the Reliability constraint but has no direct effect on the Governance constraint. This interdependence is the core reason that harness design requires simultaneous optimization across pillars, not sequential optimization of each pillar independently.
\end{remark}

The Governance Information Budget unifies the pillars by showing that they all constrain different aspects of the same fundamental problem: closing the abstraction gap reliably, efficiently, and coherently. The binding constraint rotates as the harness configuration changes: early in development (low $k$, high residual), the Abstraction and Information pillars dominate; at scale (high $k$, fast $\lambda$), the Coordination and Governance pillars dominate; for mature, stable systems (low $\lambda$, long-lived context), the Temporal pillar dominates through staleness accumulation.


% ====================================================================
% 10. EMPIRICAL FOUNDATIONS
% ====================================================================
\section{Empirical Foundations}
\label{sec:empirical}

The formal framework developed in Sections~\ref{sec:abstraction}--\ref{sec:synthesis} rests on empirical evidence of varying quality. An honest assessment of this evidence is essential.

\subsection{Strong Empirical Foundations}

\paragraph{GitClear 2025.} Analysis of 211 million changed lines across a multi-year period \citep{gitclear2025}. Key findings: code blocks with 5+ duplicated lines increased approximately 8$\times$ (overall duplication rate rose from 8.3\% to 12.3\%); refactoring collapsed from roughly 25\% to under 10\% of changed lines. Strengths: large sample, longitudinal design, objective metrics. Limitations: observational (no causal identification), single platform, code-level metrics only (no quality assessment).

\paragraph{DORA 2025.} Telemetry from over 10,000 developers \citep{dora2025}. Key finding: individual task completion increased 21\% with AI adoption, but organizational-level throughput stayed flat and software delivery stability showed a negative correlation with AI usage. Strengths: large sample, standardized metrics, multi-organization. Limitations: survey-based components, correlation not causation, no controlled intervention.

\paragraph{Kim et al.\ 2025.} 180 configurations across 5 architectures, 3 LLM families, and 4 benchmarks \citep{kim2025scaling}. Key finding: $17.2\times$ error amplification for independent agents, saturation of scaling benefits at approximately 45\% of the task completion rate. Strengths: systematic, controlled, comprehensive coverage of configuration space. Limitations: benchmark tasks (not production code), specific models and benchmarks.

\paragraph{METR 2026.} Expert evaluation of SWE-bench test-passing pull requests \citep{metr2026swebench}. Key finding: roughly 50\% of test-passing PRs would not be merged by repository maintainers (but only 68\% of human-written golden patches would be merged on re-review, so the AI-attributable gap is approximately 18 percentage points). Strengths: expert evaluation, direct relevance to the circular validation problem. Limitations: subjective judgments, small evaluator pool (4 maintainers, 3 repositories), specific benchmark context.

\subsection{Moderate Empirical Foundations}

\paragraph{Context degradation studies.} Multiple studies \citep{chroma2025contextrot, liu2024lost, du2025context} demonstrating monotonic performance degradation with context length. The power-law model $\delta(n) = (n/W_{\text{eff}})^{\alpha_d}$ with $\alpha_d \in [0.3, 0.5]$ is a fit to this data, not a derived result. The fit is reasonable but not definitive.

\paragraph{CodeRabbit 2025.} Analysis of 470 pull requests \citep{coderabbit2025}. Findings: 1.68$\times$ more issues in AI-generated PRs; 2.74$\times$ more XSS vulnerabilities. Strengths: direct comparison, automated analysis. Limitations: modest sample size, single tool, potential selection bias.

\paragraph{Hariharan et al.\ 2025.} Evaluation of plan-and-execute agent architectures \citep{hariharan2025plan}. Key finding: 63\% failure rate on 100-step tasks. Strengths: systematic scaling study. Limitations: specific agent architecture, benchmark-only evaluation.

\subsection{Weak or Missing Empirical Foundations}

Several key claims in the framework lack strong empirical support:

\begin{itemize}
    \item The optimal context budget ($|C^*| \approx 0.15W$ to $0.40W$) is derived from model fits, not from direct optimization experiments.
    \item The structural enforcement compliance gap $(1-\epsilon)^T$ is a theoretical prediction; no controlled study has measured $\epsilon$ for different enforcement mechanisms.
    \item The governance capacity bound (Theorem~\ref{thm:governance-capacity}) is a theoretical result by construction; calibrating $R_{\text{drift}}$ and $C_{\text{gov}}$ for real systems remains an open challenge.
    \item The Accretion Category is well-motivated by observational data but has not been validated through controlled experiments.
    \item VIH has not been measured in production harness settings.
    \item The cache-staleness EOQ theorem has not been empirically validated.
\end{itemize}

\subsection{Evidence Quality Assessment}
\label{sec:evidence-quality}

Table~\ref{tab:evidence-quality} provides a systematic quality assessment of the empirical foundations supporting the framework. Quality ratings follow a four-tier scheme: \textbf{High} (large sample, controlled or quasi-experimental, replicable); \textbf{Medium} (moderate sample, systematic but observational, partially replicable); \textbf{Low} (small sample, anecdotal or single-platform, limited replicability); \textbf{Very Low} (theoretical prediction only, no direct empirical measurement).

\begin{table}[H]
\centering
\caption{Evidence quality assessment for major empirical sources and theoretical claims. ``Key Gap'' identifies the most important missing piece for each data category.}
\label{tab:evidence-quality}
\small
\begin{tabular}{p{3cm}cccp{4.2cm}}
\toprule
\textbf{Data Category} & \textbf{Quality} & \textbf{Sample} & \textbf{Design} & \textbf{Key Gap} \\
\midrule
GitClear (duplication, refactoring) & High & 211M lines & Longitudinal, observational & No causal identification; single platform; no direct quality measure \\
\addlinespace
DORA 2025 (throughput, stability) & High & 10,000+ devs & Cross-sectional survey + telemetry & Survey components; no controlled intervention; correlation only \\
\addlinespace
Kim et al.\ (error amplification) & High & 180 configs & Controlled, systematic & Benchmark tasks only; 3 model families; not production code \\
\addlinespace
METR (merge gap) & Medium & $\sim$100 PRs & Expert evaluation & 4 evaluators, 3 repos; subjective criteria; small sample \\
\addlinespace
Liu et al.\ (Lost in Middle) & Medium & Multiple models & Controlled, synthetic tasks & Synthetic needle tasks; unclear production generalization \\
\addlinespace
Chroma (context rot) & Medium & 18 models & Controlled, multi-provider & Synthetic benchmarks; degradation model is a fit, not derived \\
\addlinespace
JetBrains (complexity) & Medium & Survey + metrics & Cross-sectional & Self-reported AI usage; single IDE ecosystem \\
\addlinespace
CodeRabbit (PR quality) & Medium & 470 PRs & Comparative & Single tool; potential selection bias; modest sample \\
\addlinespace
Veracode (security) & Medium & Undisclosed & Static analysis comparison & Methodology details limited; no controlled experiment \\
\addlinespace
Apollo Research (scheming) & Low & Case studies & Behavioral evaluation & Small-scale; specific models; adversarial setup \\
\addlinespace
NIST CAISI (cheating) & Low & Workshop report & Expert panel consensus & No quantitative data; policy-oriented rather than empirical \\
\addlinespace
Optimal context budget & Very Low & None & Model fit to degradation data & No direct optimization experiment; parameter sensitivity unknown \\
\addlinespace
Structural enforcement $\epsilon$ & Very Low & None & Theoretical prediction & No measurement of per-step violation probability \\
\addlinespace
VIH in production & Very Low & None & Proposed metric, unmeasured & No harness instrumentation; no baseline data \\
\addlinespace
Governance capacity ($C_{\text{gov}}$) & Very Low & None & Theoretical construction & $R_{\text{drift}}$ and $C_{\text{gov}}$ undefined operationally \\
\addlinespace
Accretion detection & Very Low & None & Category proposed, unvalidated & No labeled dataset; no precision/recall measurement \\
\addlinespace
Cache-staleness EOQ & Very Low & None & Theoretical analog & No measurement of staleness cost or optimal refresh interval \\
\bottomrule
\end{tabular}
\end{table}

\subsection{The Calibration Gap}

A recurring pattern across all seven pillars is the gap between theoretical models and empirical calibration. The models are well-grounded in established formal machinery, and the qualitative predictions align with practitioner experience. But the quantitative parameters (decay rates, threshold values, optimal counts) are either estimated from limited data or derived from model assumptions rather than measured directly.

This gap is not a criticism of the framework; it is a statement of the current state of the field. Closing this gap requires controlled experiments, which in turn requires access to production harness telemetry. Section~\ref{sec:limitations} proposes specific experiments, organized as a concrete verification roadmap with falsification criteria for each major formal result.


% ====================================================================
% 12. LIMITATIONS AND OPEN PROBLEMS
% ====================================================================
